% Generated human-review packet template.  Source records, Lean statements,
% and raw-source-to-Spec screening fields are substituted by
% scripts/review_dashboard_packet.py.
% Do not hand-edit generated packets; annotate the PDF or reconcile notes into
% the paper-local human-review log.
\documentclass[10pt]{article}
\usepackage[margin=0.43in]{geometry}
\usepackage{fontspec}
\setmainfont{Latin Modern Roman}
\setmonofont{DejaVu Sans Mono}
\usepackage{hyperref}
\usepackage{amssymb}
\usepackage{fvextra}
\usepackage{enumitem}
\setlength{\parindent}{0pt}
\setlength{\parskip}{0.15em}
\linespread{0.92}
\hypersetup{colorlinks=true,linkcolor=blue,urlcolor=blue}
\DefineVerbatimEnvironment{ReviewVerbatim}{Verbatim}{breaklines=true,breakanywhere=true,fontsize=\scriptsize}
\newcommand{\reviewerbox}{%
  \par\noindent\fbox{\parbox{0.96\linewidth}{\rule{0pt}{0.38in}}}\par
}
\newcommand{\reviewmatch}[1]{%
  \par\noindent\textbf{Reviewer check:} $\square$\; #1\par
  \noindent\emph{If left unchecked, explain the discrepancy or uncertainty below.}\par
}

\begin{document}
\fontsize{9}{10}\selectfont

\begin{center}
{\LARGE Human review packet: GKGMM19IterativeLocalVoting}\\[0.35em]
\small Generated from byte-pinned source excerpts, the expanded PaperInterface
specifications, and hash-bound raw-source-to-Spec screening records.
\end{center}

\noindent\textbf{Paper:} GKGMM19IterativeLocalVoting\\
\textbf{Paper id:} \texttt{GKGMM19IterativeLocalVoting}\\
\textbf{Source version:} not recorded\\
\textbf{Source record:} \texttt{audit/paper\_statement\_map.json}\\
\textbf{Rows:} 44\\
\textbf{Generated:} 2026-08-18\\
\textbf{Source URL:} \href{https://arxiv.org/pdf/1702.07984}{official source}

\noindent\fbox{\parbox{0.96\linewidth}{\textbf{Review status:} 0/44 source-claim screens; 0/8 library prerequisite screens. This is a review worksheet while those direct checks remain pending; it is not a final validation receipt.}}\par



\section*{How to use this packet}
For each source claim, compare the exact source excerpts with the expanded
PaperInterface specification. Mark up this PDF or fill the blank reviewer box in the TeX source;
return the annotated file for reconciliation into the paper-local human-review
log.

\section*{Open the interactive dashboard}
The interactive dashboard is an optional alternative to using this PDF; it
presents the same review surface rather than an additional review requirement.
After forking and cloning (or pulling) EconCSLib, run the following from the
repository root. The cache command is needed only the first time in a new clone.
\begin{ReviewVerbatim}
cd <your-EconCSLib-checkout>
lake exe cache get
python3 scripts/review_dashboard.py --paper GKGMM19IterativeLocalVoting --serve
\end{ReviewVerbatim}
Open \url{http://127.0.0.1:8765/} in a browser and keep that terminal running
while reviewing. The dashboard records saved annotations in this paper's
reviewer-owned review record.

\section*{Contents}
\begin{itemize}[leftmargin=1.2em]
\item \hyperlink{library-section-7af9edee67c68478}{Semantic library prerequisites (8; not paper claims)}
\begin{itemize}[leftmargin=1.2em]
\item \hyperlink{library-prerequisite-72a4fa8a1ca104da}{EconCSLib.FiniteDimensionalNorms.coordinateLinearFunctional}
\item \hyperlink{library-prerequisite-3c892f92f5d89bc4}{EconCSLib.FiniteDimensionalNorms.l1}
\item \hyperlink{library-prerequisite-9808eaf6fd3c24d8}{EconCSLib.FiniteDimensionalNorms.l2Sq}
\item \hyperlink{library-prerequisite-bef8f23ca7c4cf59}{EconCSLib.FiniteDimensionalNorms.l2}
\item \hyperlink{library-prerequisite-8d70eb5c3f7e8c76}{EconCSLib.FiniteDimensionalNorms.linf}
\item \hyperlink{library-prerequisite-8131285d8f3e27f4}{EconCSLib.FiniteDimensionalNorms.lpPower}
\item \hyperlink{library-prerequisite-c6e15042a716d8a0}{EconCSLib.FiniteDimensionalNorms.lp}
\item \hyperlink{library-prerequisite-851b58af0cd102a3}{EconCSLib.Probability.HasBoundedDensity}
\end{itemize}
\item \hyperlink{source-section-f10b5f68e4acf3dc}{Source claims (44)}
\begin{itemize}[leftmargin=1.2em]
\item \hyperlink{source-claim-db9c5a5fe969e1de}{algorithm1\_local\_neighborhood\_formulaSpec}
\item \hyperlink{source-claim-8457e0d182fe249b}{algorithm1\_norm\_projection\_formulaSpec}
\item \hyperlink{source-claim-9a151b770a276076}{algorithm1\_projected\_trajectory\_feasible\_of\_normProjectionSpec}
\item \hyperlink{source-claim-cef46a3dab6bb18f}{algorithm1\_projected\_update\_formulaSpec}
\item \hyperlink{source-claim-64105395dd6a3dd4}{algorithm1\_radius\_formulaSpec}
\item \hyperlink{source-claim-6bcb00a1cec69ec9}{algorithm1\_radius\_not\_summableSpec}
\item \hyperlink{source-claim-feefad880b5b8cc8}{algorithm1\_radius\_sq\_summableSpec}
\item \hyperlink{source-claim-0b205aab018f3bcb}{algorithm1\_radius\_ssgm\_step\_size\_conditionsSpec}
\item \hyperlink{source-claim-ce6ed26ad4d98cd3}{algorithm1\_radius\_sum\_tendsto\_atTopSpec}
\item \hyperlink{source-claim-ed1c239a4d71f093}{algorithm1\_radius\_tendsto\_zeroSpec}
\item \hyperlink{source-claim-f5582541b0eca601}{algorithm1\_stop\_condition\_formulaSpec}
\item \hyperlink{source-claim-d399110cb2008aaf}{algorithm1\_window\_stable\_formulaSpec}
\item \hyperlink{source-claim-5733b74f6b39bb81}{c3\_product\_density\_coordinate\_noncollision\_aeSpec}
\item \hyperlink{source-claim-695aae8edbfcdf18}{conditions\_c123\_formulaSpec}
\item \hyperlink{source-claim-92f9aa83d7f7fe74}{definition1\_finite\_coordinate\_l1\_formulaSpec}
\item \hyperlink{source-claim-99e53b00c8291417}{definition1\_finite\_coordinate\_l2\_formulaSpec}
\item \hyperlink{source-claim-0d61cc7c1cd21b16}{definition1\_finite\_coordinate\_linf\_formulaSpec}
\item \hyperlink{source-claim-27c15d13368b1478}{definition1\_finite\_coordinate\_lp\_formulaSpec}
\item \hyperlink{source-claim-2cc921d535fb2213}{definition1\_lp\_normed\_utilities\_formulaSpec}
\item \hyperlink{source-claim-c46295147d5004f8}{definition2\_weighted\_euclidean\_utilities\_formulaSpec}
\item \hyperlink{source-claim-a8574105e720bec5}{definition3\_decomposable\_utilities\_formulaSpec}
\item \hyperlink{source-claim-a973f59ed0e1cc66}{lemma3\_coordinate\_equality\_bad\_event\_null\_from\_boundedDensitySpec}
\item \hyperlink{source-claim-4de262cc3cbceb90}{lemma3\_coordinate\_equality\_bad\_event\_null\_from\_productMeasureSpec}
\item \hyperlink{source-claim-06e42dc4cb47855b}{lemma3\_coordinate\_noncollision\_ae\_from\_productMeasureSpec}
\item \hyperlink{source-claim-ecf761bbd6bd7ecd}{lemma3\_finite\_holder\_dual\_gradient\_candidate\_norm\_formulaSpec}
\item \hyperlink{source-claim-b6876f8d9704eca1}{lemma3\_gradient\_candidate\_hasFDerivAt\_formulaSpec}
\item \hyperlink{source-claim-086c0a9eaf7439f7}{lemma3\_gradient\_candidate\_source\_formulaSpec}
\item \hyperlink{source-claim-7c06ca676934c014}{modelA\_response\_formulaSpec}
\item \hyperlink{source-claim-2415911cc7314d54}{modelA\_response\_isMaxOn\_formulaSpec}
\item \hyperlink{source-claim-eab2a66508b56f30}{modelA\_response\_lp\_normed\_cost\_minimizer\_formulaSpec}
\item \hyperlink{source-claim-a9ebc621f1ff8c32}{modelB\_finite\_response\_formulaSpec}
\item \hyperlink{source-claim-057affe0833911ec}{modelB\_finite\_response\_neg\_lp\_cost\_gradient\_formulaSpec}
\item \hyperlink{source-claim-0bfa667fe6808e39}{proposition1\_weighted\_euclidean\_l2Spec}
\item \hyperlink{source-claim-732e45903404fd5e}{proposition2\_decomposable\_linf\_mediansSpec}
\item \hyperlink{source-claim-4f41e995c4d500c9}{appendix\_theorem4\_expected\_subgradient\_boundarySpec}
\item \hyperlink{source-claim-0f5dabcd51f758b4}{appendix\_theorem5\_ssgm\_convergence\_boundarySpec}
\item \hyperlink{source-claim-6ed3807a13e71f2b}{paper\_definition4\_dlcd\_formulaSpec}
\item \hyperlink{source-claim-9b6fb18f85b1c2c8}{dlcdBudgetUtilitySpec}
\item \hyperlink{source-claim-0c5fe9d4037e07b4}{theorem1\_lp\_normed\_dual\_casesSpec}
\item \hyperlink{source-claim-ef1266fd9eeea090}{theorem1\_norm\_pair\_l1\_linfSpec}
\item \hyperlink{source-claim-1ffdcf0343191a12}{theorem1\_norm\_pair\_l2\_l2Spec}
\item \hyperlink{source-claim-eab074ea3e4eb48e}{theorem1\_norm\_pair\_linf\_l1Spec}
\item \hyperlink{source-claim-8fdb269066d0c098}{theorem2\_modelB\_holder\_dual\_normsSpec}
\item \hyperlink{source-claim-970bf26e54daa2aa}{theorem3\_statement\_of\_full\_sampled\_projected\_source\_semantics\_univSpec}
\end{itemize}
\end{itemize}

\clearpage
\hypertarget{library-section-7af9edee67c68478}{}
\section*{Material library prerequisites}
\hypertarget{library-prerequisite-72a4fa8a1ca104da}{}
\subsection*{\small\ttfamily\raggedright EconCSLib.\allowbreak{}FiniteDimensionalNorms.\allowbreak{}coordinateLinearFunctional}
\textbf{Source connection:} no explicit byte-pinned paper source connection is registered.
\paragraph{Lean-expanded library semantic target}
\begin{ReviewVerbatim}
{ι : Type u_1} → [inst : Fintype ι] → (g : ι → ℝ) → ∑ i, g i • ContinuousLinearMap.proj i
\end{ReviewVerbatim}

\paragraph{Exact Lean library declaration}
\begin{ReviewVerbatim}
def coordinateLinearFunctional {ι : Type*} [Fintype ι]
    (g : ι → ℝ) : (ι → ℝ) →L[ℝ] ℝ :=
  ∑ i : ι, g i •
    (ContinuousLinearMap.proj (R := ℝ) (φ := fun _ : ι => ℝ) i :
      (ι → ℝ) →L[ℝ] ℝ)
\end{ReviewVerbatim}

\paragraph{Recorded source-to-library screening}
\textbf{Verdict:} \texttt{not recorded}
\reviewmatch{Matches source input}
\noindent\textbf{Reviewer annotation}\par
\reviewerbox
\clearpage
\hypertarget{library-prerequisite-3c892f92f5d89bc4}{}
\subsection*{\small\ttfamily\raggedright EconCSLib.\allowbreak{}FiniteDimensionalNorms.\allowbreak{}l1}
\textbf{Source connection:} no explicit byte-pinned paper source connection is registered.
\paragraph{Lean-expanded library semantic target}
\begin{ReviewVerbatim}
{ι : Type u_1} → [inst : Fintype ι] → (x : ι → ℝ) → ∑ i, |x i|
\end{ReviewVerbatim}

\paragraph{Exact Lean library declaration}
\begin{ReviewVerbatim}
def l1 {ι : Type*} [Fintype ι] (x : ι → ℝ) : ℝ :=
  ∑ i : ι, |x i|
\end{ReviewVerbatim}

\paragraph{Recorded source-to-library screening}
\textbf{Verdict:} \texttt{not recorded}
\reviewmatch{Matches source input}
\noindent\textbf{Reviewer annotation}\par
\reviewerbox
\clearpage
\hypertarget{library-prerequisite-9808eaf6fd3c24d8}{}
\subsection*{\small\ttfamily\raggedright EconCSLib.\allowbreak{}FiniteDimensionalNorms.\allowbreak{}l2Sq}
\textbf{Source connection:} no explicit byte-pinned paper source connection is registered.
\paragraph{Lean-expanded library semantic target}
\begin{ReviewVerbatim}
{ι : Type u_1} → [inst : Fintype ι] → (x : ι → ℝ) → ∑ i, x i ^ 2
\end{ReviewVerbatim}

\paragraph{Exact Lean library declaration}
\begin{ReviewVerbatim}
def l2Sq {ι : Type*} [Fintype ι] (x : ι → ℝ) : ℝ :=
  ∑ i : ι, x i ^ 2
\end{ReviewVerbatim}

\paragraph{Recorded source-to-library screening}
\textbf{Verdict:} \texttt{not recorded}
\reviewmatch{Matches source input}
\noindent\textbf{Reviewer annotation}\par
\reviewerbox
\clearpage
\hypertarget{library-prerequisite-bef8f23ca7c4cf59}{}
\subsection*{\small\ttfamily\raggedright EconCSLib.\allowbreak{}FiniteDimensionalNorms.\allowbreak{}l2}
\textbf{Source connection:} no explicit byte-pinned paper source connection is registered.
\paragraph{Lean-expanded library semantic target}
\begin{ReviewVerbatim}
{ι : Type u_1} → [inst : Fintype ι] → (x : ι → ℝ) → √(EconCSLib.FiniteDimensionalNorms.l2Sq x)
\end{ReviewVerbatim}

\paragraph{Exact Lean library declaration}
\begin{ReviewVerbatim}
def l2 {ι : Type*} [Fintype ι] (x : ι → ℝ) : ℝ :=
  Real.sqrt (l2Sq x)
\end{ReviewVerbatim}

\paragraph{Recorded source-to-library screening}
\textbf{Verdict:} \texttt{not recorded}
\reviewmatch{Matches source input}
\noindent\textbf{Reviewer annotation}\par
\reviewerbox
\clearpage
\hypertarget{library-prerequisite-8d70eb5c3f7e8c76}{}
\subsection*{\small\ttfamily\raggedright EconCSLib.\allowbreak{}FiniteDimensionalNorms.\allowbreak{}linf}
\textbf{Source connection:} no explicit byte-pinned paper source connection is registered.
\paragraph{Lean-expanded library semantic target}
\begin{ReviewVerbatim}
{ι : Type u_1} → [inst : Fintype ι] → [inst_1 : Nonempty ι] → (x : ι → ℝ) → Finset.univ.sup' ⋯ fun i => |x i|
\end{ReviewVerbatim}

\paragraph{Exact Lean library declaration}
\begin{ReviewVerbatim}
def linf {ι : Type*} [Fintype ι] [Nonempty ι] (x : ι → ℝ) : ℝ :=
  (Finset.univ : Finset ι).sup' Finset.univ_nonempty (fun i => |x i|)
\end{ReviewVerbatim}

\paragraph{Recorded source-to-library screening}
\textbf{Verdict:} \texttt{not recorded}
\reviewmatch{Matches source input}
\noindent\textbf{Reviewer annotation}\par
\reviewerbox
\clearpage
\hypertarget{library-prerequisite-8131285d8f3e27f4}{}
\subsection*{\small\ttfamily\raggedright EconCSLib.\allowbreak{}FiniteDimensionalNorms.\allowbreak{}lpPower}
\textbf{Source connection:} no explicit byte-pinned paper source connection is registered.
\paragraph{Lean-expanded library semantic target}
\begin{ReviewVerbatim}
{ι : Type u_1} → [inst : Fintype ι] → (p : ℝ) → (x : ι → ℝ) → ∑ i, |x i| ^ p
\end{ReviewVerbatim}

\paragraph{Exact Lean library declaration}
\begin{ReviewVerbatim}
def lpPower {ι : Type*} [Fintype ι] (p : ℝ) (x : ι → ℝ) : ℝ :=
  ∑ i : ι, |x i| ^ p
\end{ReviewVerbatim}

\paragraph{Recorded source-to-library screening}
\textbf{Verdict:} \texttt{not recorded}
\reviewmatch{Matches source input}
\noindent\textbf{Reviewer annotation}\par
\reviewerbox
\clearpage
\hypertarget{library-prerequisite-c6e15042a716d8a0}{}
\subsection*{\small\ttfamily\raggedright EconCSLib.\allowbreak{}FiniteDimensionalNorms.\allowbreak{}lp}
\textbf{Source connection:} no explicit byte-pinned paper source connection is registered.
\paragraph{Lean-expanded library semantic target}
\begin{ReviewVerbatim}
{ι : Type u_1} → [inst : Fintype ι] → (p : ℝ) → (x : ι → ℝ) → EconCSLib.FiniteDimensionalNorms.lpPower p x ^ (1 / p)
\end{ReviewVerbatim}

\paragraph{Exact Lean library declaration}
\begin{ReviewVerbatim}
def lp {ι : Type*} [Fintype ι] (p : ℝ) (x : ι → ℝ) : ℝ :=
  (lpPower p x) ^ (1 / p)
\end{ReviewVerbatim}

\paragraph{Recorded source-to-library screening}
\textbf{Verdict:} \texttt{not recorded}
\reviewmatch{Matches source input}
\noindent\textbf{Reviewer annotation}\par
\reviewerbox
\clearpage
\hypertarget{library-prerequisite-851b58af0cd102a3}{}
\subsection*{\small\ttfamily\raggedright EconCSLib.\allowbreak{}Probability.\allowbreak{}HasBoundedDensity}
\textbf{Source connection:} no explicit byte-pinned paper source connection is registered.
\paragraph{Lean-expanded library semantic target}
\begin{ReviewVerbatim}
∀ {α : Type u_1} [inst : MeasurableSpace α] (ν μ : MeasureTheory.Measure α) (C : ENNReal),
  ∃ D, μ = ν.withDensity D ∧ ∀ᵐ (x : α) ∂ν, D x ≤ C
\end{ReviewVerbatim}

\paragraph{Exact Lean library declaration}
\begin{ReviewVerbatim}
def HasBoundedDensity
    {α : Type*} [MeasurableSpace α]
    (ν μ : Measure α) (C : ℝ≥0∞) : Prop :=
  ∃ D : α → ℝ≥0∞, μ = ν.withDensity D ∧ ∀ᵐ x ∂ν, D x ≤ C
\end{ReviewVerbatim}

\paragraph{Recorded source-to-library screening}
\textbf{Verdict:} \texttt{not recorded}
\reviewmatch{Matches source input}
\noindent\textbf{Reviewer annotation}\par
\reviewerbox

\clearpage
\hypertarget{source-claim-db9c5a5fe969e1de}{}
\hypertarget{source-section-f10b5f68e4acf3dc}{}
\section*{Source claims}
\subsection*{Review row 1}
\textbf{Source locator:} {\footnotesize\raggedright cited publication, lines 170-\allowbreak{}192\par}
\paragraph{Verbatim source input}
\begin{ReviewVerbatim}
Algorithm 1: Iterative Local Voting (ILV)

Inputs:
Initial solution x0 ∈X, tolerance ϵ > 0, an integer N, initial radius r0 > 0,
termination time T, norm q for local neighborhood.
Output:
Solution x.

• For t ≥1, sample a voter vt ∈V at random from the population; set rt = r0/t
and elicit

x′
t = arg maxx∈{s:||s−xt−1||q≤rt}fvt(x),
(1)

and then compute xt = [x′
t]X , where [·]X is a projection onto space X; i.e. ask
the voter to move to her favorite point within constraints, and then project the
reported point onto X.

• Stop when either t = T, in which case return xT , or when
maxl, m∈{t−N,...,t} |xl −xm| ≤ϵ, in which case return x = xt.
\end{ReviewVerbatim}


\paragraph{Expanded PaperInterface specification}
\begin{ReviewVerbatim}
∀ {Point : Type u_1} (G : GKGMM19IterativeLocalVoting.ILVGeometryFormulaData Point)
  (q : GKGMM19IterativeLocalVoting.SourceNorm) (center candidate : Point) (r : ℝ),
  candidate ∈ GKGMM19IterativeLocalVoting.localNeighborhoodFormulaData G q center r ↔
    candidate ∈ G.solutionSpace ∧ G.normDistance q candidate center ≤ r
\end{ReviewVerbatim}

\paragraph{Lean proof endpoint (built separately)}\begin{ReviewVerbatim}
GKGMM19IterativeLocalVoting.algorithm1_local_neighborhood_formula
\end{ReviewVerbatim}

\paragraph{Recorded source-to-Spec screening}
\textbf{Verdict:} \texttt{not recorded}
\\
\textbf{Reason:} not recorded
\reviewmatch{Matches source input}
\noindent\textbf{Reviewer annotation}\par
\reviewerbox
\clearpage
\hypertarget{source-claim-8457e0d182fe249b}{}
\section*{Review row 2}
\textbf{Source locator:} {\footnotesize\raggedright cited publication, lines 170-\allowbreak{}192\par}
\paragraph{Verbatim source input}
\begin{ReviewVerbatim}
Algorithm 1: Iterative Local Voting (ILV)

Inputs:
Initial solution x0 ∈X, tolerance ϵ > 0, an integer N, initial radius r0 > 0,
termination time T, norm q for local neighborhood.
Output:
Solution x.

• For t ≥1, sample a voter vt ∈V at random from the population; set rt = r0/t
and elicit

x′
t = arg maxx∈{s:||s−xt−1||q≤rt}fvt(x),
(1)

and then compute xt = [x′
t]X , where [·]X is a projection onto space X; i.e. ask
the voter to move to her favorite point within constraints, and then project the
reported point onto X.

• Stop when either t = T, in which case return xT , or when
maxl, m∈{t−N,...,t} |xl −xm| ≤ϵ, in which case return x = xt.
\end{ReviewVerbatim}


\paragraph{Expanded PaperInterface specification}
\begin{ReviewVerbatim}
∀ {Point : Type u_1} (G : GKGMM19IterativeLocalVoting.ILVGeometryFormulaData Point)
  (q : GKGMM19IterativeLocalVoting.SourceNorm) (project : Point → Point),
  GKGMM19IterativeLocalVoting.isNormProjectionOntoFormulaData G q project ↔
    ∀ (y : Point), project y ∈ G.solutionSpace ∧ IsMinOn (fun x => G.normDistance q x y) G.solutionSpace (project y)
\end{ReviewVerbatim}

\paragraph{Lean proof endpoint (built separately)}\begin{ReviewVerbatim}
GKGMM19IterativeLocalVoting.algorithm1_norm_projection_formula
\end{ReviewVerbatim}

\paragraph{Recorded source-to-Spec screening}
\textbf{Verdict:} \texttt{not recorded}
\\
\textbf{Reason:} not recorded
\reviewmatch{Matches source input}
\noindent\textbf{Reviewer annotation}\par
\reviewerbox
\clearpage
\hypertarget{source-claim-9a151b770a276076}{}
\section*{Review row 3}
\textbf{Source locator:} {\footnotesize\raggedright cited publication, lines 170-\allowbreak{}192\par}
\paragraph{Verbatim source input}
\begin{ReviewVerbatim}
Algorithm 1: Iterative Local Voting (ILV)

Inputs:
Initial solution x0 ∈X, tolerance ϵ > 0, an integer N, initial radius r0 > 0,
termination time T, norm q for local neighborhood.
Output:
Solution x.

• For t ≥1, sample a voter vt ∈V at random from the population; set rt = r0/t
and elicit

x′
t = arg maxx∈{s:||s−xt−1||q≤rt}fvt(x),
(1)

and then compute xt = [x′
t]X , where [·]X is a projection onto space X; i.e. ask
the voter to move to her favorite point within constraints, and then project the
reported point onto X.

• Stop when either t = T, in which case return xT , or when
maxl, m∈{t−N,...,t} |xl −xm| ≤ϵ, in which case return x = xt.
\end{ReviewVerbatim}


\paragraph{Expanded PaperInterface specification}
\begin{ReviewVerbatim}
∀ {Point : Type u_1} {G : GKGMM19IterativeLocalVoting.ILVGeometryFormulaData Point}
  {q : GKGMM19IterativeLocalVoting.SourceNorm} {project : Point → Point} {raw trajectory : ℕ → Point},
  GKGMM19IterativeLocalVoting.isNormProjectionOntoFormulaData G q project →
    trajectory 0 ∈ G.solutionSpace →
      (∀ (t : ℕ), GKGMM19IterativeLocalVoting.Algorithm1ProjectedUpdate project (raw t) (trajectory (t + 1))) →
        ∀ (t : ℕ), trajectory t ∈ G.solutionSpace
\end{ReviewVerbatim}

\paragraph{Lean proof endpoint (built separately)}\begin{ReviewVerbatim}
GKGMM19IterativeLocalVoting.algorithm1_projected_trajectory_feasible_of_normProjection
\end{ReviewVerbatim}

\paragraph{Recorded source-to-Spec screening}
\textbf{Verdict:} \texttt{not recorded}
\\
\textbf{Reason:} not recorded
\reviewmatch{Matches source input}
\noindent\textbf{Reviewer annotation}\par
\reviewerbox
\clearpage
\hypertarget{source-claim-cef46a3dab6bb18f}{}
\section*{Review row 4}
\textbf{Source locator:} {\footnotesize\raggedright cited publication, lines 170-\allowbreak{}192\par}
\paragraph{Verbatim source input}
\begin{ReviewVerbatim}
Algorithm 1: Iterative Local Voting (ILV)

Inputs:
Initial solution x0 ∈X, tolerance ϵ > 0, an integer N, initial radius r0 > 0,
termination time T, norm q for local neighborhood.
Output:
Solution x.

• For t ≥1, sample a voter vt ∈V at random from the population; set rt = r0/t
and elicit

x′
t = arg maxx∈{s:||s−xt−1||q≤rt}fvt(x),
(1)

and then compute xt = [x′
t]X , where [·]X is a projection onto space X; i.e. ask
the voter to move to her favorite point within constraints, and then project the
reported point onto X.

• Stop when either t = T, in which case return xT , or when
maxl, m∈{t−N,...,t} |xl −xm| ≤ϵ, in which case return x = xt.
\end{ReviewVerbatim}


\paragraph{Expanded PaperInterface specification}
\begin{ReviewVerbatim}
∀ {Point : Type u_1} (project : Point → Point) (raw next : Point),
  GKGMM19IterativeLocalVoting.Algorithm1ProjectedUpdate project raw next ↔ next = project raw
\end{ReviewVerbatim}

\paragraph{Lean proof endpoint (built separately)}\begin{ReviewVerbatim}
GKGMM19IterativeLocalVoting.algorithm1_projected_update_formula
\end{ReviewVerbatim}

\paragraph{Recorded source-to-Spec screening}
\textbf{Verdict:} \texttt{not recorded}
\\
\textbf{Reason:} not recorded
\reviewmatch{Matches source input}
\noindent\textbf{Reviewer annotation}\par
\reviewerbox
\clearpage
\hypertarget{source-claim-64105395dd6a3dd4}{}
\section*{Review row 5}
\textbf{Source locator:} {\footnotesize\raggedright cited publication, lines 170-\allowbreak{}192\par}
\paragraph{Verbatim source input}
\begin{ReviewVerbatim}
Algorithm 1: Iterative Local Voting (ILV)

Inputs:
Initial solution x0 ∈X, tolerance ϵ > 0, an integer N, initial radius r0 > 0,
termination time T, norm q for local neighborhood.
Output:
Solution x.

• For t ≥1, sample a voter vt ∈V at random from the population; set rt = r0/t
and elicit

x′
t = arg maxx∈{s:||s−xt−1||q≤rt}fvt(x),
(1)

and then compute xt = [x′
t]X , where [·]X is a projection onto space X; i.e. ask
the voter to move to her favorite point within constraints, and then project the
reported point onto X.

• Stop when either t = T, in which case return xT , or when
maxl, m∈{t−N,...,t} |xl −xm| ≤ϵ, in which case return x = xt.
\end{ReviewVerbatim}


\paragraph{Expanded PaperInterface specification}
\begin{ReviewVerbatim}
∀ (r0 : ℝ) (t : ℕ), GKGMM19IterativeLocalVoting.ilvRadius r0 t = r0 / ↑t
\end{ReviewVerbatim}

\paragraph{Lean proof endpoint (built separately)}\begin{ReviewVerbatim}
GKGMM19IterativeLocalVoting.algorithm1_radius_formula
\end{ReviewVerbatim}

\paragraph{Recorded source-to-Spec screening}
\textbf{Verdict:} \texttt{not recorded}
\\
\textbf{Reason:} not recorded
\reviewmatch{Matches source input}
\noindent\textbf{Reviewer annotation}\par
\reviewerbox
\clearpage
\hypertarget{source-claim-6bcb00a1cec69ec9}{}
\section*{Review row 6}
\textbf{Source locator:} {\footnotesize\raggedright cited publication, lines 170-\allowbreak{}192\par}
\paragraph{Verbatim source input}
\begin{ReviewVerbatim}
Algorithm 1: Iterative Local Voting (ILV)

Inputs:
Initial solution x0 ∈X, tolerance ϵ > 0, an integer N, initial radius r0 > 0,
termination time T, norm q for local neighborhood.
Output:
Solution x.

• For t ≥1, sample a voter vt ∈V at random from the population; set rt = r0/t
and elicit

x′
t = arg maxx∈{s:||s−xt−1||q≤rt}fvt(x),
(1)

and then compute xt = [x′
t]X , where [·]X is a projection onto space X; i.e. ask
the voter to move to her favorite point within constraints, and then project the
reported point onto X.

• Stop when either t = T, in which case return xT , or when
maxl, m∈{t−N,...,t} |xl −xm| ≤ϵ, in which case return x = xt.
\end{ReviewVerbatim}


\paragraph{Expanded PaperInterface specification}
\begin{ReviewVerbatim}
∀ {r0 : ℝ}, 0 < r0 → ¬Summable fun t => GKGMM19IterativeLocalVoting.ilvRadius r0 (t + 1)
\end{ReviewVerbatim}

\paragraph{Lean proof endpoint (built separately)}\begin{ReviewVerbatim}
GKGMM19IterativeLocalVoting.algorithm1_radius_not_summable
\end{ReviewVerbatim}

\paragraph{Recorded source-to-Spec screening}
\textbf{Verdict:} \texttt{not recorded}
\\
\textbf{Reason:} not recorded
\reviewmatch{Matches source input}
\noindent\textbf{Reviewer annotation}\par
\reviewerbox
\clearpage
\hypertarget{source-claim-feefad880b5b8cc8}{}
\section*{Review row 7}
\textbf{Source locator:} {\footnotesize\raggedright cited publication, lines 170-\allowbreak{}192\par}
\paragraph{Verbatim source input}
\begin{ReviewVerbatim}
Algorithm 1: Iterative Local Voting (ILV)

Inputs:
Initial solution x0 ∈X, tolerance ϵ > 0, an integer N, initial radius r0 > 0,
termination time T, norm q for local neighborhood.
Output:
Solution x.

• For t ≥1, sample a voter vt ∈V at random from the population; set rt = r0/t
and elicit

x′
t = arg maxx∈{s:||s−xt−1||q≤rt}fvt(x),
(1)

and then compute xt = [x′
t]X , where [·]X is a projection onto space X; i.e. ask
the voter to move to her favorite point within constraints, and then project the
reported point onto X.

• Stop when either t = T, in which case return xT , or when
maxl, m∈{t−N,...,t} |xl −xm| ≤ϵ, in which case return x = xt.
\end{ReviewVerbatim}


\paragraph{Expanded PaperInterface specification}
\begin{ReviewVerbatim}
∀ (r0 : ℝ), Summable fun t => GKGMM19IterativeLocalVoting.ilvRadius r0 (t + 1) ^ 2
\end{ReviewVerbatim}

\paragraph{Lean proof endpoint (built separately)}\begin{ReviewVerbatim}
GKGMM19IterativeLocalVoting.algorithm1_radius_sq_summable
\end{ReviewVerbatim}

\paragraph{Recorded source-to-Spec screening}
\textbf{Verdict:} \texttt{not recorded}
\\
\textbf{Reason:} not recorded
\reviewmatch{Matches source input}
\noindent\textbf{Reviewer annotation}\par
\reviewerbox
\clearpage
\hypertarget{source-claim-0b205aab018f3bcb}{}
\section*{Review row 8}
\textbf{Source locator:} {\footnotesize\raggedright cited publication, lines 170-\allowbreak{}192\par}
\paragraph{Verbatim source input}
\begin{ReviewVerbatim}
Algorithm 1: Iterative Local Voting (ILV)

Inputs:
Initial solution x0 ∈X, tolerance ϵ > 0, an integer N, initial radius r0 > 0,
termination time T, norm q for local neighborhood.
Output:
Solution x.

• For t ≥1, sample a voter vt ∈V at random from the population; set rt = r0/t
and elicit

x′
t = arg maxx∈{s:||s−xt−1||q≤rt}fvt(x),
(1)

and then compute xt = [x′
t]X , where [·]X is a projection onto space X; i.e. ask
the voter to move to her favorite point within constraints, and then project the
reported point onto X.

• Stop when either t = T, in which case return xT , or when
maxl, m∈{t−N,...,t} |xl −xm| ≤ϵ, in which case return x = xt.
\end{ReviewVerbatim}


\paragraph{Expanded PaperInterface specification}
\begin{ReviewVerbatim}
∀ {r0 : ℝ}, 0 < r0 → GKGMM19IterativeLocalVoting.SSGMStepSizeConditions (GKGMM19IterativeLocalVoting.ilvRadius r0)
\end{ReviewVerbatim}

\paragraph{Lean proof endpoint (built separately)}\begin{ReviewVerbatim}
GKGMM19IterativeLocalVoting.algorithm1_radius_ssgm_step_size_conditions
\end{ReviewVerbatim}

\paragraph{Recorded source-to-Spec screening}
\textbf{Verdict:} \texttt{not recorded}
\\
\textbf{Reason:} not recorded
\reviewmatch{Matches source input}
\noindent\textbf{Reviewer annotation}\par
\reviewerbox
\clearpage
\hypertarget{source-claim-ce6ed26ad4d98cd3}{}
\section*{Review row 9}
\textbf{Source locator:} {\footnotesize\raggedright cited publication, lines 170-\allowbreak{}192\par}
\paragraph{Verbatim source input}
\begin{ReviewVerbatim}
Algorithm 1: Iterative Local Voting (ILV)

Inputs:
Initial solution x0 ∈X, tolerance ϵ > 0, an integer N, initial radius r0 > 0,
termination time T, norm q for local neighborhood.
Output:
Solution x.

• For t ≥1, sample a voter vt ∈V at random from the population; set rt = r0/t
and elicit

x′
t = arg maxx∈{s:||s−xt−1||q≤rt}fvt(x),
(1)

and then compute xt = [x′
t]X , where [·]X is a projection onto space X; i.e. ask
the voter to move to her favorite point within constraints, and then project the
reported point onto X.

• Stop when either t = T, in which case return xT , or when
maxl, m∈{t−N,...,t} |xl −xm| ≤ϵ, in which case return x = xt.
\end{ReviewVerbatim}


\paragraph{Expanded PaperInterface specification}
\begin{ReviewVerbatim}
∀ {r0 : ℝ},
  0 < r0 →
    Filter.Tendsto (fun n => ∑ t ∈ Finset.range n, GKGMM19IterativeLocalVoting.ilvRadius r0 (t + 1)) Filter.atTop
      Filter.atTop
\end{ReviewVerbatim}

\paragraph{Lean proof endpoint (built separately)}\begin{ReviewVerbatim}
GKGMM19IterativeLocalVoting.algorithm1_radius_sum_tendsto_atTop
\end{ReviewVerbatim}

\paragraph{Recorded source-to-Spec screening}
\textbf{Verdict:} \texttt{not recorded}
\\
\textbf{Reason:} not recorded
\reviewmatch{Matches source input}
\noindent\textbf{Reviewer annotation}\par
\reviewerbox
\clearpage
\hypertarget{source-claim-ed1c239a4d71f093}{}
\section*{Review row 10}
\textbf{Source locator:} {\footnotesize\raggedright cited publication, lines 170-\allowbreak{}192\par}
\paragraph{Verbatim source input}
\begin{ReviewVerbatim}
Algorithm 1: Iterative Local Voting (ILV)

Inputs:
Initial solution x0 ∈X, tolerance ϵ > 0, an integer N, initial radius r0 > 0,
termination time T, norm q for local neighborhood.
Output:
Solution x.

• For t ≥1, sample a voter vt ∈V at random from the population; set rt = r0/t
and elicit

x′
t = arg maxx∈{s:||s−xt−1||q≤rt}fvt(x),
(1)

and then compute xt = [x′
t]X , where [·]X is a projection onto space X; i.e. ask
the voter to move to her favorite point within constraints, and then project the
reported point onto X.

• Stop when either t = T, in which case return xT , or when
maxl, m∈{t−N,...,t} |xl −xm| ≤ϵ, in which case return x = xt.
\end{ReviewVerbatim}


\paragraph{Expanded PaperInterface specification}
\begin{ReviewVerbatim}
∀ (r0 : ℝ), Filter.Tendsto (GKGMM19IterativeLocalVoting.ilvRadius r0) Filter.atTop (nhds 0)
\end{ReviewVerbatim}

\paragraph{Lean proof endpoint (built separately)}\begin{ReviewVerbatim}
GKGMM19IterativeLocalVoting.algorithm1_radius_tendsto_zero
\end{ReviewVerbatim}

\paragraph{Recorded source-to-Spec screening}
\textbf{Verdict:} \texttt{not recorded}
\\
\textbf{Reason:} not recorded
\reviewmatch{Matches source input}
\noindent\textbf{Reviewer annotation}\par
\reviewerbox
\clearpage
\hypertarget{source-claim-f5582541b0eca601}{}
\section*{Review row 11}
\textbf{Source locator:} {\footnotesize\raggedright cited publication, lines 170-\allowbreak{}192\par}
\paragraph{Verbatim source input}
\begin{ReviewVerbatim}
Algorithm 1: Iterative Local Voting (ILV)

Inputs:
Initial solution x0 ∈X, tolerance ϵ > 0, an integer N, initial radius r0 > 0,
termination time T, norm q for local neighborhood.
Output:
Solution x.

• For t ≥1, sample a voter vt ∈V at random from the population; set rt = r0/t
and elicit

x′
t = arg maxx∈{s:||s−xt−1||q≤rt}fvt(x),
(1)

and then compute xt = [x′
t]X , where [·]X is a projection onto space X; i.e. ask
the voter to move to her favorite point within constraints, and then project the
reported point onto X.

• Stop when either t = T, in which case return xT , or when
maxl, m∈{t−N,...,t} |xl −xm| ≤ϵ, in which case return x = xt.
\end{ReviewVerbatim}


\paragraph{Expanded PaperInterface specification}
\begin{ReviewVerbatim}
∀ {Point : Type u_1} (G : GKGMM19IterativeLocalVoting.ILVGeometryFormulaData Point)
  (q : GKGMM19IterativeLocalVoting.SourceNorm) (trajectory : ℕ → Point) (T N t : ℕ) (epsilon : ℝ),
  GKGMM19IterativeLocalVoting.algorithm1StopConditionFormulaData G q trajectory T N t epsilon ↔
    t = T ∨ GKGMM19IterativeLocalVoting.algorithm1WindowStableFormulaData G q trajectory t N epsilon
\end{ReviewVerbatim}

\paragraph{Lean proof endpoint (built separately)}\begin{ReviewVerbatim}
GKGMM19IterativeLocalVoting.algorithm1_stop_condition_formula
\end{ReviewVerbatim}

\paragraph{Recorded source-to-Spec screening}
\textbf{Verdict:} \texttt{not recorded}
\\
\textbf{Reason:} not recorded
\reviewmatch{Matches source input}
\noindent\textbf{Reviewer annotation}\par
\reviewerbox
\clearpage
\hypertarget{source-claim-d399110cb2008aaf}{}
\section*{Review row 12}
\textbf{Source locator:} {\footnotesize\raggedright cited publication, lines 170-\allowbreak{}192\par}
\paragraph{Verbatim source input}
\begin{ReviewVerbatim}
Algorithm 1: Iterative Local Voting (ILV)

Inputs:
Initial solution x0 ∈X, tolerance ϵ > 0, an integer N, initial radius r0 > 0,
termination time T, norm q for local neighborhood.
Output:
Solution x.

• For t ≥1, sample a voter vt ∈V at random from the population; set rt = r0/t
and elicit

x′
t = arg maxx∈{s:||s−xt−1||q≤rt}fvt(x),
(1)

and then compute xt = [x′
t]X , where [·]X is a projection onto space X; i.e. ask
the voter to move to her favorite point within constraints, and then project the
reported point onto X.

• Stop when either t = T, in which case return xT , or when
maxl, m∈{t−N,...,t} |xl −xm| ≤ϵ, in which case return x = xt.
\end{ReviewVerbatim}


\paragraph{Expanded PaperInterface specification}
\begin{ReviewVerbatim}
∀ {Point : Type u_1} (G : GKGMM19IterativeLocalVoting.ILVGeometryFormulaData Point)
  (q : GKGMM19IterativeLocalVoting.SourceNorm) (trajectory : ℕ → Point) (t N : ℕ) (epsilon : ℝ),
  GKGMM19IterativeLocalVoting.algorithm1WindowStableFormulaData G q trajectory t N epsilon ↔
    ∀ (l m : ℕ),
      l ∈ Finset.Icc (t - N) t → m ∈ Finset.Icc (t - N) t → G.normDistance q (trajectory l) (trajectory m) ≤ epsilon
\end{ReviewVerbatim}

\paragraph{Lean proof endpoint (built separately)}\begin{ReviewVerbatim}
GKGMM19IterativeLocalVoting.algorithm1_window_stable_formula
\end{ReviewVerbatim}

\paragraph{Recorded source-to-Spec screening}
\textbf{Verdict:} \texttt{not recorded}
\\
\textbf{Reason:} not recorded
\reviewmatch{Matches source input}
\noindent\textbf{Reviewer annotation}\par
\reviewerbox
\clearpage
\hypertarget{source-claim-5733b74f6b39bb81}{}
\section*{Review row 13}
\textbf{Source locator:} {\footnotesize\raggedright cited publication, lines 2138-\allowbreak{}2155\par}
\paragraph{Verbatim source input}
\begin{ReviewVerbatim}
Lemma 3. ∀(p, q) s.t. p > 0, q > 0, and 1/p + 1/q = 1, ||∇||x −xv||p||q = 1, ∀x s.t.
xm ̸= xm
v for any m.

Theorem 2. Suppose that conditions C1, C2, and C3 are satis[U+001C control character]ed, the voter utilities are
Lp normed, and voters respond to query (1) according to Model B. Then, ILV with Lq

neighborhoods converges to the societal optimal point w.p. 1 for any p > 0 and q > 0 such
that 1/p + 1/q = 1.

Proof. Since the probability of picking a voter v such that xm
t = xm
v for some dimension m
is 0, we have ˜gvt(xt) = gvt for gvt = ∇fvt(xt). Thus we obtain the gradient exactly, and
hence Theorem 5 applies with bt = 0 for all t.

C.5 Proof of Propositions
\end{ReviewVerbatim}


\paragraph{Expanded PaperInterface specification}
\begin{ReviewVerbatim}
∀ {Coord : Type u_1} [inst : Fintype Coord]
  (D : GKGMM19IterativeLocalVoting.FiniteCoordinateIdealDistributionData Coord) (x : Coord → ℝ),
  ∀ᵐ (ideal : Coord → ℝ) ∂D.idealMeasure, ∀ (i : Coord), x i ≠ ideal i
\end{ReviewVerbatim}

\paragraph{Lean proof endpoint (built separately)}\begin{ReviewVerbatim}
GKGMM19IterativeLocalVoting.c3_product_density_coordinate_noncollision_ae
\end{ReviewVerbatim}

\paragraph{Recorded source-to-Spec screening}
\textbf{Verdict:} \texttt{not recorded}
\\
\textbf{Reason:} not recorded
\reviewmatch{Matches source input}
\noindent\textbf{Reviewer annotation}\par
\reviewerbox
\clearpage
\hypertarget{source-claim-695aae8edbfcdf18}{}
\section*{Review row 14}
\textbf{Source locator:} {\footnotesize\raggedright cited publication, lines 410-\allowbreak{}420\par}
\paragraph{Verbatim source input}
\begin{ReviewVerbatim}
behavior models. For the rest of the technical analysis, we make the following assumptions
on our model.

C1 The solution space X ⊆RM is non-empty, bounded, closed, and convex.

C2 Each voter v has a unique ideal solution xv ∈X.

C3 The ideal point xv of each voter is drawn independently from a probability distribution
with a bounded and measurable density function hX .

Under this model, for a solution x ∈X, the societal utility is given by Ev[fv(x)]. and the
\end{ReviewVerbatim}


\paragraph{Expanded PaperInterface specification}
\begin{ReviewVerbatim}
∀
  (solutionSpace_nonempty_bounded_closed_convex uniqueIdealSolutions idealDistribution_bounded_measurable_density :
    Prop),
  solutionSpace_nonempty_bounded_closed_convex ∧ uniqueIdealSolutions ∧ idealDistribution_bounded_measurable_density ↔
    solutionSpace_nonempty_bounded_closed_convex ∧ uniqueIdealSolutions ∧ idealDistribution_bounded_measurable_density
\end{ReviewVerbatim}

\paragraph{Lean proof endpoint (built separately)}\begin{ReviewVerbatim}
GKGMM19IterativeLocalVoting.conditions_c123_formula
\end{ReviewVerbatim}

\paragraph{Recorded source-to-Spec screening}
\textbf{Verdict:} \texttt{not recorded}
\\
\textbf{Reason:} not recorded
\reviewmatch{Matches source input}
\noindent\textbf{Reviewer annotation}\par
\reviewerbox
\clearpage
\hypertarget{source-claim-92f9aa83d7f7fe74}{}
\section*{Review row 15}
\textbf{Source locator:} {\footnotesize\raggedright cited publication, lines 430-\allowbreak{}450\par}
\paragraph{Verbatim source input}
\begin{ReviewVerbatim}
3.1 Spatial Utilities

Here we consider spatial utility functions, where the utilities of each voters can be expressed
in the form of some kind of spatial distance from their ideal solutions. First, we consider
the following kind of utilities.

De[U+001C control character]nition 1. Lp normed utilities. The voter utility function is Lp normed if fv(x) =
−||x −xv||p, ∀x ∈X.

Under such utilities, for p = 1, 2 and ∞, restricting voters to a ball in the dual norm
leads to convergence to the societal optimum.

Theorem 1. Suppose that conditions C1, C2, and C3 are satis[U+001C control character]ed, the voter utilities are Lp

normed, and voters respond to query (1) according to either Model A or Model B. Then,
ILV with Lq neighborhoods converges to the societal optimal point w.p. 1 when (p, q) = (2, 2),
(1, ∞), or (∞, 1).

322
\end{ReviewVerbatim}


\paragraph{Expanded PaperInterface specification}
\begin{ReviewVerbatim}
∀ {Voter : Type u_1} {Coord : Type u_2} [inst : Fintype Coord] [inst_1 : Nonempty Coord]
  (utility : Voter → (Coord → ℝ) → ℝ) (ideal : Voter → Coord → ℝ),
  (∀ (v : Voter) (x : Coord → ℝ),
      utility v x =
        -GKGMM19IterativeLocalVoting.finiteCoordinateDistance GKGMM19IterativeLocalVoting.SourceNorm.l1 x (ideal v)) ↔
    ∀ (v : Voter) (x : Coord → ℝ), utility v x = -EconCSLib.FiniteDimensionalNorms.l1 fun m => x m - ideal v m
\end{ReviewVerbatim}

\paragraph{Lean proof endpoint (built separately)}\begin{ReviewVerbatim}
GKGMM19IterativeLocalVoting.definition1_finite_coordinate_l1_formula
\end{ReviewVerbatim}

\paragraph{Recorded source-to-Spec screening}
\textbf{Verdict:} \texttt{not recorded}
\\
\textbf{Reason:} not recorded
\reviewmatch{Matches source input}
\noindent\textbf{Reviewer annotation}\par
\reviewerbox
\clearpage
\hypertarget{source-claim-99e53b00c8291417}{}
\section*{Review row 16}
\textbf{Source locator:} {\footnotesize\raggedright cited publication, lines 430-\allowbreak{}450\par}
\paragraph{Verbatim source input}
\begin{ReviewVerbatim}
3.1 Spatial Utilities

Here we consider spatial utility functions, where the utilities of each voters can be expressed
in the form of some kind of spatial distance from their ideal solutions. First, we consider
the following kind of utilities.

De[U+001C control character]nition 1. Lp normed utilities. The voter utility function is Lp normed if fv(x) =
−||x −xv||p, ∀x ∈X.

Under such utilities, for p = 1, 2 and ∞, restricting voters to a ball in the dual norm
leads to convergence to the societal optimum.

Theorem 1. Suppose that conditions C1, C2, and C3 are satis[U+001C control character]ed, the voter utilities are Lp

normed, and voters respond to query (1) according to either Model A or Model B. Then,
ILV with Lq neighborhoods converges to the societal optimal point w.p. 1 when (p, q) = (2, 2),
(1, ∞), or (∞, 1).

322
\end{ReviewVerbatim}


\paragraph{Expanded PaperInterface specification}
\begin{ReviewVerbatim}
∀ {Voter : Type u_1} {Coord : Type u_2} [inst : Fintype Coord] [inst_1 : Nonempty Coord]
  (utility : Voter → (Coord → ℝ) → ℝ) (ideal : Voter → Coord → ℝ),
  (∀ (v : Voter) (x : Coord → ℝ),
      utility v x =
        -GKGMM19IterativeLocalVoting.finiteCoordinateDistance GKGMM19IterativeLocalVoting.SourceNorm.l2 x (ideal v)) ↔
    ∀ (v : Voter) (x : Coord → ℝ), utility v x = -EconCSLib.FiniteDimensionalNorms.l2 fun m => x m - ideal v m
\end{ReviewVerbatim}

\paragraph{Lean proof endpoint (built separately)}\begin{ReviewVerbatim}
GKGMM19IterativeLocalVoting.definition1_finite_coordinate_l2_formula
\end{ReviewVerbatim}

\paragraph{Recorded source-to-Spec screening}
\textbf{Verdict:} \texttt{not recorded}
\\
\textbf{Reason:} not recorded
\reviewmatch{Matches source input}
\noindent\textbf{Reviewer annotation}\par
\reviewerbox
\clearpage
\hypertarget{source-claim-0d61cc7c1cd21b16}{}
\section*{Review row 17}
\textbf{Source locator:} {\footnotesize\raggedright cited publication, lines 430-\allowbreak{}450\par}
\paragraph{Verbatim source input}
\begin{ReviewVerbatim}
3.1 Spatial Utilities

Here we consider spatial utility functions, where the utilities of each voters can be expressed
in the form of some kind of spatial distance from their ideal solutions. First, we consider
the following kind of utilities.

De[U+001C control character]nition 1. Lp normed utilities. The voter utility function is Lp normed if fv(x) =
−||x −xv||p, ∀x ∈X.

Under such utilities, for p = 1, 2 and ∞, restricting voters to a ball in the dual norm
leads to convergence to the societal optimum.

Theorem 1. Suppose that conditions C1, C2, and C3 are satis[U+001C control character]ed, the voter utilities are Lp

normed, and voters respond to query (1) according to either Model A or Model B. Then,
ILV with Lq neighborhoods converges to the societal optimal point w.p. 1 when (p, q) = (2, 2),
(1, ∞), or (∞, 1).

322
\end{ReviewVerbatim}


\paragraph{Expanded PaperInterface specification}
\begin{ReviewVerbatim}
∀ {Voter : Type u_1} {Coord : Type u_2} [inst : Fintype Coord] [inst_1 : Nonempty Coord]
  (utility : Voter → (Coord → ℝ) → ℝ) (ideal : Voter → Coord → ℝ),
  (∀ (v : Voter) (x : Coord → ℝ),
      utility v x =
        -GKGMM19IterativeLocalVoting.finiteCoordinateDistance GKGMM19IterativeLocalVoting.SourceNorm.linfty x
            (ideal v)) ↔
    ∀ (v : Voter) (x : Coord → ℝ), utility v x = -EconCSLib.FiniteDimensionalNorms.linf fun m => x m - ideal v m
\end{ReviewVerbatim}

\paragraph{Lean proof endpoint (built separately)}\begin{ReviewVerbatim}
GKGMM19IterativeLocalVoting.definition1_finite_coordinate_linf_formula
\end{ReviewVerbatim}

\paragraph{Recorded source-to-Spec screening}
\textbf{Verdict:} \texttt{not recorded}
\\
\textbf{Reason:} not recorded
\reviewmatch{Matches source input}
\noindent\textbf{Reviewer annotation}\par
\reviewerbox
\clearpage
\hypertarget{source-claim-27c15d13368b1478}{}
\section*{Review row 18}
\textbf{Source locator:} {\footnotesize\raggedright cited publication, lines 430-\allowbreak{}450\par}
\paragraph{Verbatim source input}
\begin{ReviewVerbatim}
3.1 Spatial Utilities

Here we consider spatial utility functions, where the utilities of each voters can be expressed
in the form of some kind of spatial distance from their ideal solutions. First, we consider
the following kind of utilities.

De[U+001C control character]nition 1. Lp normed utilities. The voter utility function is Lp normed if fv(x) =
−||x −xv||p, ∀x ∈X.

Under such utilities, for p = 1, 2 and ∞, restricting voters to a ball in the dual norm
leads to convergence to the societal optimum.

Theorem 1. Suppose that conditions C1, C2, and C3 are satis[U+001C control character]ed, the voter utilities are Lp

normed, and voters respond to query (1) according to either Model A or Model B. Then,
ILV with Lq neighborhoods converges to the societal optimal point w.p. 1 when (p, q) = (2, 2),
(1, ∞), or (∞, 1).

322
\end{ReviewVerbatim}


\paragraph{Expanded PaperInterface specification}
\begin{ReviewVerbatim}
∀ {Voter : Type u_1} {Coord : Type u_2} [inst : Fintype Coord] [inst_1 : Nonempty Coord]
  (utility : Voter → (Coord → ℝ) → ℝ) (ideal : Voter → Coord → ℝ) (p : ℝ),
  (∀ (v : Voter) (x : Coord → ℝ),
      utility v x =
        -GKGMM19IterativeLocalVoting.finiteCoordinateDistance (GKGMM19IterativeLocalVoting.SourceNorm.lp p) x
            (ideal v)) ↔
    ∀ (v : Voter) (x : Coord → ℝ), utility v x = -EconCSLib.FiniteDimensionalNorms.lp p fun m => x m - ideal v m
\end{ReviewVerbatim}

\paragraph{Lean proof endpoint (built separately)}\begin{ReviewVerbatim}
GKGMM19IterativeLocalVoting.definition1_finite_coordinate_lp_formula
\end{ReviewVerbatim}

\paragraph{Recorded source-to-Spec screening}
\textbf{Verdict:} \texttt{not recorded}
\\
\textbf{Reason:} not recorded
\reviewmatch{Matches source input}
\noindent\textbf{Reviewer annotation}\par
\reviewerbox
\clearpage
\hypertarget{source-claim-2cc921d535fb2213}{}
\section*{Review row 19}
\textbf{Source locator:} {\footnotesize\raggedright cited publication, lines 430-\allowbreak{}450\par}
\paragraph{Verbatim source input}
\begin{ReviewVerbatim}
3.1 Spatial Utilities

Here we consider spatial utility functions, where the utilities of each voters can be expressed
in the form of some kind of spatial distance from their ideal solutions. First, we consider
the following kind of utilities.

De[U+001C control character]nition 1. Lp normed utilities. The voter utility function is Lp normed if fv(x) =
−||x −xv||p, ∀x ∈X.

Under such utilities, for p = 1, 2 and ∞, restricting voters to a ball in the dual norm
leads to convergence to the societal optimum.

Theorem 1. Suppose that conditions C1, C2, and C3 are satis[U+001C control character]ed, the voter utilities are Lp

normed, and voters respond to query (1) according to either Model A or Model B. Then,
ILV with Lq neighborhoods converges to the societal optimal point w.p. 1 when (p, q) = (2, 2),
(1, ∞), or (∞, 1).

322
\end{ReviewVerbatim}


\paragraph{Expanded PaperInterface specification}
\begin{ReviewVerbatim}
∀ {Voter : Type u_1} {Point : Type u_2} (U : GKGMM19IterativeLocalVoting.ILVUtilityFormulaData Voter Point)
  (ideal : Voter → Point) (p : GKGMM19IterativeLocalVoting.SourceNorm),
  GKGMM19IterativeLocalVoting.isLpNormedUtilitiesFormulaData U ideal p ↔
    ∀ (v : Voter) (x : Point), U.utility v x = -U.normDistance p x (ideal v)
\end{ReviewVerbatim}

\paragraph{Lean proof endpoint (built separately)}\begin{ReviewVerbatim}
GKGMM19IterativeLocalVoting.definition1_lp_normed_utilities_formula
\end{ReviewVerbatim}

\paragraph{Recorded source-to-Spec screening}
\textbf{Verdict:} \texttt{not recorded}
\\
\textbf{Reason:} not recorded
\reviewmatch{Matches source input}
\noindent\textbf{Reviewer annotation}\par
\reviewerbox
\clearpage
\hypertarget{source-claim-c46295147d5004f8}{}
\section*{Review row 20}
\textbf{Source locator:} {\footnotesize\raggedright cited publication, lines 480-\allowbreak{}535\par}
\paragraph{Verbatim source input}
\begin{ReviewVerbatim}
sampling such a voter decreases fast enough).
Next, we introduce another general class of utility functions, which we call Weighted
Euclidean utilities, for which one can obtain convergence to a unique solution.

De[U+001C control character]nition 2. Weighted Euclidean utilities. Let the solution space X be decomposable
into K di[U+001B control character]erent sub-spaces, so that x = (x1, . . . , xK) for each x ∈X (where �K
k=1 dim(xk) =
M). Suppose that the utility function of the voter v is

fv(x) = −

K
�

k=1

wk
v
||wv||2
||xk −xk
v||2.

where wv is a voter-speci[U+001C control character]c weight vector, then the function is a Weighted Euclidean utility
function. We further assume that wv ∈W ⊂RK
+ and xv are independently drawn for each
voter v from a joint probability distribution with a bounded and measurable density function,
with W nonempty, bounded, closed, and convex.

This utility function can be interpreted as follows: the decision-making problem is de-
composable into K sub-problems, and each voter v has an ideal point xk
v and a weight wk
v
for each sub-problem k, so that the voter's disutility for a solution is the weighted sum of
the Euclidean distances to the ideal points in each sub-problems. Such utility functions
may emerge in facility location problems, for example, where voters have preferences on

323


[form-feed in source extraction]

Garg, Kamble, Goel, Marn, & Munagala

the locations of multiple facilities on a map. This utility form is also the one most closely
related to the existing literature on Directional Equilibria and Quadratic Voting, in which
preferences are linear. To recover the weighted linear preferences case, set K = M, with
each sub-space of dimension 1. In this case, the following holds:

Proposition 1. Suppose that conditions C1, C2, and C3 are satis[U+001C control character]ed, the voter utilities are
Weighted Euclidean, and voters correctly respond to query (1) according to either Model A
or Model B. Then, ILV with L2 neighborhoods converges with probability 1 to the societal
optimal point.

The intuition for the result is as follows: as long as the neighborhood does not contain
the ideal point of the sampled voter, the correct response to query (1) under weighted
Euclidean preferences is to move the solution in the direction of the ideal point to the
neighborhood boundary, which, as it turns out, is the same as the direction of the gradient.
Thus with radius rt, the e[U+001B control character]ective movement is
\end{ReviewVerbatim}


\paragraph{Expanded PaperInterface specification}
\begin{ReviewVerbatim}
∀ {Voter : Type u_1} {Point : Type u_2} {Component : Type u_3} (utility : Voter → Point → ℝ)
  (W : GKGMM19IterativeLocalVoting.WeightedEuclideanStructure Voter Point Component),
  (W.weightsAndIdealsDistributionCondition ∧
      ∀ (v : Voter) (x : Point),
        utility v x = -∑ k ∈ W.components, W.weight v k / W.weightNorm2 v * W.componentDistance k x v) ↔
    W.weightsAndIdealsDistributionCondition ∧
      ∀ (v : Voter) (x : Point),
        utility v x = -∑ k ∈ W.components, W.weight v k / W.weightNorm2 v * W.componentDistance k x v
\end{ReviewVerbatim}

\paragraph{Lean proof endpoint (built separately)}\begin{ReviewVerbatim}
GKGMM19IterativeLocalVoting.definition2_weighted_euclidean_utilities_formula
\end{ReviewVerbatim}

\paragraph{Recorded source-to-Spec screening}
\textbf{Verdict:} \texttt{not recorded}
\\
\textbf{Reason:} not recorded
\reviewmatch{Matches source input}
\noindent\textbf{Reviewer annotation}\par
\reviewerbox
\clearpage
\hypertarget{source-claim-a8574105e720bec5}{}
\section*{Review row 21}
\textbf{Source locator:} {\footnotesize\raggedright cited publication, lines 540-\allowbreak{}575\par}
\paragraph{Verbatim source input}
\begin{ReviewVerbatim}
normalized (i.e. not divided by ||w||2), the algorithm would converge to the same point, as
if utility functions were normalized.

3.2 Decomposable Utilities

Next consider the general class of decomposable utilities, motivated by the fact that the
algorithm with L∞neighborhoods is of special interest since they are easy for humans to
understand: one can change each dimension up to a certain amount, independent of the
others.

De[U+001C control character]nition 3. Decomposable utilities. A voter utility function is decomposable if there
exists concave functions fm
v
for m ∈{1 . . . M} such that fv(x) = �M
m=1 fm
v (xm).

If the utility functions for the voters are decomposable, then we can show that our
algorithm under L∞neighborhoods converges to the vector of medians of voters' ideal points
on each dimension. Suppose that hm
X is the marginal density function of the random variable
xm
v , and let ¯xm be the set of medians of xm
v . (By set of medians, we mean the set of points
such that, on each dimension, the mass of voters with ideal points above and below.)

Proposition 2. Suppose that conditions C1, C2, and C3 are satis[U+001C control character]ed, the voter utilities
are decomposable, and voters respond to query (1) according to either Model A or Model
B. Then, ILV with L∞neighborhoods converges with probability 1 to a point in the set of
medians ¯x.

Although simply eliciting each agent's optimal solution and computing the vector of
median allocations on each dimension is a viable approach in the case of decomposable
utilities, deciding an optimal allocation across multiple dimensions is a more challenging
cognitive task than deciding whether one wants to increase or decrease each dimension
relative to the current solution (see Section 5.2.3 for experimental evidence). In fact, in this
\end{ReviewVerbatim}


\paragraph{Expanded PaperInterface specification}
\begin{ReviewVerbatim}
∀ {Voter : Type u_1} {Point : Type u_2} {Coord : Type u_3} (utility : Voter → Point → ℝ)
  (D : GKGMM19IterativeLocalVoting.DecomposableStructure Voter Point Coord),
  (D.coordinateUtilitiesConcave ∧
      ∀ (v : Voter) (x : Point), utility v x = ∑ m ∈ D.coords, D.coordinateUtility m v (D.coordinate m x)) ↔
    D.coordinateUtilitiesConcave ∧
      ∀ (v : Voter) (x : Point), utility v x = ∑ m ∈ D.coords, D.coordinateUtility m v (D.coordinate m x)
\end{ReviewVerbatim}

\paragraph{Lean proof endpoint (built separately)}\begin{ReviewVerbatim}
GKGMM19IterativeLocalVoting.definition3_decomposable_utilities_formula
\end{ReviewVerbatim}

\paragraph{Recorded source-to-Spec screening}
\textbf{Verdict:} \texttt{not recorded}
\\
\textbf{Reason:} not recorded
\reviewmatch{Matches source input}
\noindent\textbf{Reviewer annotation}\par
\reviewerbox
\clearpage
\hypertarget{source-claim-a973f59ed0e1cc66}{}
\section*{Review row 22}
\textbf{Source locator:} {\footnotesize\raggedright cited publication, lines 2138-\allowbreak{}2155\par}
\paragraph{Verbatim source input}
\begin{ReviewVerbatim}
Lemma 3. ∀(p, q) s.t. p > 0, q > 0, and 1/p + 1/q = 1, ||∇||x −xv||p||q = 1, ∀x s.t.
xm ̸= xm
v for any m.

Theorem 2. Suppose that conditions C1, C2, and C3 are satis[U+001C control character]ed, the voter utilities are
Lp normed, and voters respond to query (1) according to Model B. Then, ILV with Lq

neighborhoods converges to the societal optimal point w.p. 1 for any p > 0 and q > 0 such
that 1/p + 1/q = 1.

Proof. Since the probability of picking a voter v such that xm
t = xm
v for some dimension m
is 0, we have ˜gvt(xt) = gvt for gvt = ∇fvt(xt). Thus we obtain the gradient exactly, and
hence Theorem 5 applies with bt = 0 for all t.

C.5 Proof of Propositions
\end{ReviewVerbatim}


\paragraph{Expanded PaperInterface specification}
\begin{ReviewVerbatim}
∀ {Coord : Type u_1} [Fintype Coord] [inst : MeasurableSpace (Coord → ℝ)] {ν μ : MeasureTheory.Measure (Coord → ℝ)}
  {C : ENNReal},
  EconCSLib.Probability.HasBoundedDensity ν μ C →
    ∀ (x : Coord → ℝ),
      (∀ (i : Coord), ν (GKGMM19IterativeLocalVoting.coordinateEqualityHyperplane x i) = 0) →
        μ (GKGMM19IterativeLocalVoting.coordinateEqualityBadEvent x) = 0
\end{ReviewVerbatim}

\paragraph{Lean proof endpoint (built separately)}\begin{ReviewVerbatim}
GKGMM19IterativeLocalVoting.lemma3_coordinate_equality_bad_event_null_from_boundedDensity
\end{ReviewVerbatim}

\paragraph{Recorded source-to-Spec screening}
\textbf{Verdict:} \texttt{not recorded}
\\
\textbf{Reason:} not recorded
\reviewmatch{Matches source input}
\noindent\textbf{Reviewer annotation}\par
\reviewerbox
\clearpage
\hypertarget{source-claim-4de262cc3cbceb90}{}
\section*{Review row 23}
\textbf{Source locator:} {\footnotesize\raggedright cited publication, lines 2138-\allowbreak{}2155\par}
\paragraph{Verbatim source input}
\begin{ReviewVerbatim}
Lemma 3. ∀(p, q) s.t. p > 0, q > 0, and 1/p + 1/q = 1, ||∇||x −xv||p||q = 1, ∀x s.t.
xm ̸= xm
v for any m.

Theorem 2. Suppose that conditions C1, C2, and C3 are satis[U+001C control character]ed, the voter utilities are
Lp normed, and voters respond to query (1) according to Model B. Then, ILV with Lq

neighborhoods converges to the societal optimal point w.p. 1 for any p > 0 and q > 0 such
that 1/p + 1/q = 1.

Proof. Since the probability of picking a voter v such that xm
t = xm
v for some dimension m
is 0, we have ˜gvt(xt) = gvt for gvt = ∇fvt(xt). Thus we obtain the gradient exactly, and
hence Theorem 5 applies with bt = 0 for all t.

C.5 Proof of Propositions
\end{ReviewVerbatim}


\paragraph{Expanded PaperInterface specification}
\begin{ReviewVerbatim}
∀ {Coord : Type u_1} [inst : Fintype Coord] (ρ : MeasureTheory.Measure ℝ) [MeasureTheory.SigmaFinite ρ]
  [MeasureTheory.NoAtoms ρ] {μ : MeasureTheory.Measure (Coord → ℝ)} {C : ENNReal},
  EconCSLib.Probability.HasBoundedDensity (MeasureTheory.Measure.pi fun x => ρ) μ C →
    ∀ (x : Coord → ℝ), μ (GKGMM19IterativeLocalVoting.coordinateEqualityBadEvent x) = 0
\end{ReviewVerbatim}

\paragraph{Lean proof endpoint (built separately)}\begin{ReviewVerbatim}
GKGMM19IterativeLocalVoting.lemma3_coordinate_equality_bad_event_null_from_productMeasure
\end{ReviewVerbatim}

\paragraph{Recorded source-to-Spec screening}
\textbf{Verdict:} \texttt{not recorded}
\\
\textbf{Reason:} not recorded
\reviewmatch{Matches source input}
\noindent\textbf{Reviewer annotation}\par
\reviewerbox
\clearpage
\hypertarget{source-claim-06e42dc4cb47855b}{}
\section*{Review row 24}
\textbf{Source locator:} {\footnotesize\raggedright cited publication, lines 2138-\allowbreak{}2155\par}
\paragraph{Verbatim source input}
\begin{ReviewVerbatim}
Lemma 3. ∀(p, q) s.t. p > 0, q > 0, and 1/p + 1/q = 1, ||∇||x −xv||p||q = 1, ∀x s.t.
xm ̸= xm
v for any m.

Theorem 2. Suppose that conditions C1, C2, and C3 are satis[U+001C control character]ed, the voter utilities are
Lp normed, and voters respond to query (1) according to Model B. Then, ILV with Lq

neighborhoods converges to the societal optimal point w.p. 1 for any p > 0 and q > 0 such
that 1/p + 1/q = 1.

Proof. Since the probability of picking a voter v such that xm
t = xm
v for some dimension m
is 0, we have ˜gvt(xt) = gvt for gvt = ∇fvt(xt). Thus we obtain the gradient exactly, and
hence Theorem 5 applies with bt = 0 for all t.

C.5 Proof of Propositions
\end{ReviewVerbatim}


\paragraph{Expanded PaperInterface specification}
\begin{ReviewVerbatim}
∀ {Coord : Type u_1} [inst : Fintype Coord] (ρ : MeasureTheory.Measure ℝ) [MeasureTheory.SigmaFinite ρ]
  [MeasureTheory.NoAtoms ρ] {μ : MeasureTheory.Measure (Coord → ℝ)} {C : ENNReal},
  EconCSLib.Probability.HasBoundedDensity (MeasureTheory.Measure.pi fun x => ρ) μ C →
    ∀ (x : Coord → ℝ), ∀ᵐ (ideal : Coord → ℝ) ∂μ, ∀ (i : Coord), x i ≠ ideal i
\end{ReviewVerbatim}

\paragraph{Lean proof endpoint (built separately)}\begin{ReviewVerbatim}
GKGMM19IterativeLocalVoting.lemma3_coordinate_noncollision_ae_from_productMeasure
\end{ReviewVerbatim}

\paragraph{Recorded source-to-Spec screening}
\textbf{Verdict:} \texttt{not recorded}
\\
\textbf{Reason:} not recorded
\reviewmatch{Matches source input}
\noindent\textbf{Reviewer annotation}\par
\reviewerbox
\clearpage
\hypertarget{source-claim-ecf761bbd6bd7ecd}{}
\section*{Review row 25}
\textbf{Source locator:} {\footnotesize\raggedright cited publication, lines 2138-\allowbreak{}2155\par}
\paragraph{Verbatim source input}
\begin{ReviewVerbatim}
Lemma 3. ∀(p, q) s.t. p > 0, q > 0, and 1/p + 1/q = 1, ||∇||x −xv||p||q = 1, ∀x s.t.
xm ̸= xm
v for any m.

Theorem 2. Suppose that conditions C1, C2, and C3 are satis[U+001C control character]ed, the voter utilities are
Lp normed, and voters respond to query (1) according to Model B. Then, ILV with Lq

neighborhoods converges to the societal optimal point w.p. 1 for any p > 0 and q > 0 such
that 1/p + 1/q = 1.

Proof. Since the probability of picking a voter v such that xm
t = xm
v for some dimension m
is 0, we have ˜gvt(xt) = gvt for gvt = ∇fvt(xt). Thus we obtain the gradient exactly, and
hence Theorem 5 applies with bt = 0 for all t.

C.5 Proof of Propositions
\end{ReviewVerbatim}


\paragraph{Expanded PaperInterface specification}
\begin{ReviewVerbatim}
∀ {Coord : Type u_1} [inst : Fintype Coord] [inst_1 : Nonempty Coord] {p q : ℝ},
  GKGMM19IterativeLocalVoting.HolderDualFinite p q →
    ∀ {x ideal : Coord → ℝ},
      (∀ (i : Coord), x i ≠ ideal i) →
        GKGMM19IterativeLocalVoting.finiteCoordinateNorm (GKGMM19IterativeLocalVoting.SourceNorm.lp q)
            (GKGMM19IterativeLocalVoting.lpCostGradientCandidate p fun i => x i - ideal i) =
          1
\end{ReviewVerbatim}

\paragraph{Lean proof endpoint (built separately)}\begin{ReviewVerbatim}
GKGMM19IterativeLocalVoting.lemma3_finite_holder_dual_gradient_candidate_norm_formula
\end{ReviewVerbatim}

\paragraph{Recorded source-to-Spec screening}
\textbf{Verdict:} \texttt{not recorded}
\\
\textbf{Reason:} not recorded
\reviewmatch{Matches source input}
\noindent\textbf{Reviewer annotation}\par
\reviewerbox
\clearpage
\hypertarget{source-claim-b6876f8d9704eca1}{}
\section*{Review row 26}
\textbf{Source locator:} {\footnotesize\raggedright cited publication, lines 2138-\allowbreak{}2155\par}
\paragraph{Verbatim source input}
\begin{ReviewVerbatim}
Lemma 3. ∀(p, q) s.t. p > 0, q > 0, and 1/p + 1/q = 1, ||∇||x −xv||p||q = 1, ∀x s.t.
xm ̸= xm
v for any m.

Theorem 2. Suppose that conditions C1, C2, and C3 are satis[U+001C control character]ed, the voter utilities are
Lp normed, and voters respond to query (1) according to Model B. Then, ILV with Lq

neighborhoods converges to the societal optimal point w.p. 1 for any p > 0 and q > 0 such
that 1/p + 1/q = 1.

Proof. Since the probability of picking a voter v such that xm
t = xm
v for some dimension m
is 0, we have ˜gvt(xt) = gvt for gvt = ∇fvt(xt). Thus we obtain the gradient exactly, and
hence Theorem 5 applies with bt = 0 for all t.

C.5 Proof of Propositions
\end{ReviewVerbatim}


\paragraph{Expanded PaperInterface specification}
\begin{ReviewVerbatim}
∀ {Coord : Type u_1} [inst : Fintype Coord] [Nonempty Coord] {p : ℝ},
  1 < p →
    ∀ {x ideal : Coord → ℝ},
      (∀ (i : Coord), x i ≠ ideal i) →
        HasFDerivAt (fun y => EconCSLib.FiniteDimensionalNorms.lp p fun i => y i - ideal i)
          (EconCSLib.FiniteDimensionalNorms.coordinateLinearFunctional
            (GKGMM19IterativeLocalVoting.lpCostGradientCandidate p fun i => x i - ideal i))
          x
\end{ReviewVerbatim}

\paragraph{Lean proof endpoint (built separately)}\begin{ReviewVerbatim}
GKGMM19IterativeLocalVoting.lemma3_gradient_candidate_hasFDerivAt_formula
\end{ReviewVerbatim}

\paragraph{Recorded source-to-Spec screening}
\textbf{Verdict:} \texttt{not recorded}
\\
\textbf{Reason:} not recorded
\reviewmatch{Matches source input}
\noindent\textbf{Reviewer annotation}\par
\reviewerbox
\clearpage
\hypertarget{source-claim-086c0a9eaf7439f7}{}
\section*{Review row 27}
\textbf{Source locator:} {\footnotesize\raggedright cited publication, lines 2138-\allowbreak{}2155\par}
\paragraph{Verbatim source input}
\begin{ReviewVerbatim}
Lemma 3. ∀(p, q) s.t. p > 0, q > 0, and 1/p + 1/q = 1, ||∇||x −xv||p||q = 1, ∀x s.t.
xm ̸= xm
v for any m.

Theorem 2. Suppose that conditions C1, C2, and C3 are satis[U+001C control character]ed, the voter utilities are
Lp normed, and voters respond to query (1) according to Model B. Then, ILV with Lq

neighborhoods converges to the societal optimal point w.p. 1 for any p > 0 and q > 0 such
that 1/p + 1/q = 1.

Proof. Since the probability of picking a voter v such that xm
t = xm
v for some dimension m
is 0, we have ˜gvt(xt) = gvt for gvt = ∇fvt(xt). Thus we obtain the gradient exactly, and
hence Theorem 5 applies with bt = 0 for all t.

C.5 Proof of Propositions
\end{ReviewVerbatim}


\paragraph{Expanded PaperInterface specification}
\begin{ReviewVerbatim}
∀ {Coord : Type u_1} [inst : Fintype Coord] {p : ℝ},
  0 < p →
    ∀ (d : Coord → ℝ),
      GKGMM19IterativeLocalVoting.lpCostGradientCandidate p d = fun i =>
        |d i| ^ (p - 1) * (d i / |d i|) / EconCSLib.FiniteDimensionalNorms.lp p d ^ (p - 1)
\end{ReviewVerbatim}

\paragraph{Lean proof endpoint (built separately)}\begin{ReviewVerbatim}
GKGMM19IterativeLocalVoting.lemma3_gradient_candidate_source_formula
\end{ReviewVerbatim}

\paragraph{Recorded source-to-Spec screening}
\textbf{Verdict:} \texttt{not recorded}
\\
\textbf{Reason:} not recorded
\reviewmatch{Matches source input}
\noindent\textbf{Reviewer annotation}\par
\reviewerbox
\clearpage
\hypertarget{source-claim-7c06ca676934c014}{}
\section*{Review row 28}
\textbf{Source locator:} {\footnotesize\raggedright cited publication, lines 118-\allowbreak{}145\par}
\paragraph{Verbatim source input}
\begin{ReviewVerbatim}
bounds on the sum of absolute values of the changes, the sum of the square of the changes,
and the maximum change, respectively.)
More formally, consider a M-dimensional societal decision problem in X ⊂RM and a
population of voters V, where each voter v ∈V has bounded utility fv(x) ∈R, ∀x ∈X.
Then we consider the class of algorithms described in Algorithm 1. We study the algorithm
class under two plausible models of how voters respond to query (1), which asks for the
voter's favorite point in a local region.

• Model A:
One possibility is that voters exactly perform the maximization asked of
them, responding with their favorite point in the given Lq norm constraint set. In other
words, they return a point arg maxx∈{s:||s−xt−1||q≤rt}fvt(x). Note that by de[U+001C control character]nition of this
movement, the algorithm is myopically incentive compatible: if a voter is the last voter
and no projections are used, then truthfully performing this movement is the dominant
strategy. In general, the mechanism is not globally incentive compatible, nor incentive
compatible with projections onto the feasible region. Simple examples of manipulations
in both instances exist.

• Model B: On the other hand, voters may not actually search within the constraint set
to [U+001C control character]nd their favorite point inside of it. Rather, a voter v may have an idea about how
to best improve the current point and then move in that direction to the boundary of
the given constraint set. This model leads to a voter moving the current solution in the
direction of the gradient of her utility function, returning a point xt−1 + rt
gt
||gt||q , for some
gt ∈∂fvt(xt−1). Note that ∂f(x) denotes the set of subgradients of a function f at x, i.e.
g ∈∂f(x) if ∀y, f(y) −f(x) ≥gT (y −x).
\end{ReviewVerbatim}


\paragraph{Expanded PaperInterface specification}
\begin{ReviewVerbatim}
∀ {Voter : Type u_1} {Point : Type u_2} (U : GKGMM19IterativeLocalVoting.ILVUtilityFormulaData Voter Point)
  (q : GKGMM19IterativeLocalVoting.SourceNorm) (center : Point) (r : ℝ) (voter : Voter) (response : Point),
  GKGMM19IterativeLocalVoting.modelAResponseFormulaData U q center r voter response ↔
    response ∈ GKGMM19IterativeLocalVoting.localNeighborhoodFormulaData U.geometry q center r ∧
      ∀ candidate ∈ GKGMM19IterativeLocalVoting.localNeighborhoodFormulaData U.geometry q center r,
        U.utility voter candidate ≤ U.utility voter response
\end{ReviewVerbatim}

\paragraph{Lean proof endpoint (built separately)}\begin{ReviewVerbatim}
GKGMM19IterativeLocalVoting.modelA_response_formula
\end{ReviewVerbatim}

\paragraph{Recorded source-to-Spec screening}
\textbf{Verdict:} \texttt{not recorded}
\\
\textbf{Reason:} not recorded
\reviewmatch{Matches source input}
\noindent\textbf{Reviewer annotation}\par
\reviewerbox
\clearpage
\hypertarget{source-claim-2415911cc7314d54}{}
\section*{Review row 29}
\textbf{Source locator:} {\footnotesize\raggedright cited publication, lines 118-\allowbreak{}145\par}
\paragraph{Verbatim source input}
\begin{ReviewVerbatim}
bounds on the sum of absolute values of the changes, the sum of the square of the changes,
and the maximum change, respectively.)
More formally, consider a M-dimensional societal decision problem in X ⊂RM and a
population of voters V, where each voter v ∈V has bounded utility fv(x) ∈R, ∀x ∈X.
Then we consider the class of algorithms described in Algorithm 1. We study the algorithm
class under two plausible models of how voters respond to query (1), which asks for the
voter's favorite point in a local region.

• Model A:
One possibility is that voters exactly perform the maximization asked of
them, responding with their favorite point in the given Lq norm constraint set. In other
words, they return a point arg maxx∈{s:||s−xt−1||q≤rt}fvt(x). Note that by de[U+001C control character]nition of this
movement, the algorithm is myopically incentive compatible: if a voter is the last voter
and no projections are used, then truthfully performing this movement is the dominant
strategy. In general, the mechanism is not globally incentive compatible, nor incentive
compatible with projections onto the feasible region. Simple examples of manipulations
in both instances exist.

• Model B: On the other hand, voters may not actually search within the constraint set
to [U+001C control character]nd their favorite point inside of it. Rather, a voter v may have an idea about how
to best improve the current point and then move in that direction to the boundary of
the given constraint set. This model leads to a voter moving the current solution in the
direction of the gradient of her utility function, returning a point xt−1 + rt
gt
||gt||q , for some
gt ∈∂fvt(xt−1). Note that ∂f(x) denotes the set of subgradients of a function f at x, i.e.
g ∈∂f(x) if ∀y, f(y) −f(x) ≥gT (y −x).
\end{ReviewVerbatim}


\paragraph{Expanded PaperInterface specification}
\begin{ReviewVerbatim}
∀ {Voter : Type u_1} {Point : Type u_2} (U : GKGMM19IterativeLocalVoting.ILVUtilityFormulaData Voter Point)
  (q : GKGMM19IterativeLocalVoting.SourceNorm) (center : Point) (r : ℝ) (voter : Voter) (response : Point),
  GKGMM19IterativeLocalVoting.modelAResponseFormulaData U q center r voter response ↔
    response ∈ GKGMM19IterativeLocalVoting.localNeighborhoodFormulaData U.geometry q center r ∧
      IsMaxOn (U.utility voter) (GKGMM19IterativeLocalVoting.localNeighborhoodFormulaData U.geometry q center r)
        response
\end{ReviewVerbatim}

\paragraph{Lean proof endpoint (built separately)}\begin{ReviewVerbatim}
GKGMM19IterativeLocalVoting.modelA_response_isMaxOn_formula
\end{ReviewVerbatim}

\paragraph{Recorded source-to-Spec screening}
\textbf{Verdict:} \texttt{not recorded}
\\
\textbf{Reason:} not recorded
\reviewmatch{Matches source input}
\noindent\textbf{Reviewer annotation}\par
\reviewerbox
\clearpage
\hypertarget{source-claim-eab2a66508b56f30}{}
\section*{Review row 30}
\textbf{Source locator:} {\footnotesize\raggedright cited publication, lines 118-\allowbreak{}145\par}
\paragraph{Verbatim source input}
\begin{ReviewVerbatim}
bounds on the sum of absolute values of the changes, the sum of the square of the changes,
and the maximum change, respectively.)
More formally, consider a M-dimensional societal decision problem in X ⊂RM and a
population of voters V, where each voter v ∈V has bounded utility fv(x) ∈R, ∀x ∈X.
Then we consider the class of algorithms described in Algorithm 1. We study the algorithm
class under two plausible models of how voters respond to query (1), which asks for the
voter's favorite point in a local region.

• Model A:
One possibility is that voters exactly perform the maximization asked of
them, responding with their favorite point in the given Lq norm constraint set. In other
words, they return a point arg maxx∈{s:||s−xt−1||q≤rt}fvt(x). Note that by de[U+001C control character]nition of this
movement, the algorithm is myopically incentive compatible: if a voter is the last voter
and no projections are used, then truthfully performing this movement is the dominant
strategy. In general, the mechanism is not globally incentive compatible, nor incentive
compatible with projections onto the feasible region. Simple examples of manipulations
in both instances exist.

• Model B: On the other hand, voters may not actually search within the constraint set
to [U+001C control character]nd their favorite point inside of it. Rather, a voter v may have an idea about how
to best improve the current point and then move in that direction to the boundary of
the given constraint set. This model leads to a voter moving the current solution in the
direction of the gradient of her utility function, returning a point xt−1 + rt
gt
||gt||q , for some
gt ∈∂fvt(xt−1). Note that ∂f(x) denotes the set of subgradients of a function f at x, i.e.
g ∈∂f(x) if ∀y, f(y) −f(x) ≥gT (y −x).
\end{ReviewVerbatim}


\paragraph{Expanded PaperInterface specification}
\begin{ReviewVerbatim}
∀ {Voter : Type u_1} {Point : Type u_2} (U : GKGMM19IterativeLocalVoting.ILVUtilityFormulaData Voter Point)
  (p q : GKGMM19IterativeLocalVoting.SourceNorm) (ideal : Voter → Point) (center : Point) (r : ℝ) (voter : Voter)
  (response : Point),
  GKGMM19IterativeLocalVoting.isLpNormedUtilitiesFormulaData U ideal p →
    (GKGMM19IterativeLocalVoting.modelAResponseFormulaData U q center r voter response ↔
      response ∈ GKGMM19IterativeLocalVoting.localNeighborhoodFormulaData U.geometry q center r ∧
        IsMinOn (fun candidate => U.normDistance p candidate (ideal voter))
          (GKGMM19IterativeLocalVoting.localNeighborhoodFormulaData U.geometry q center r) response)
\end{ReviewVerbatim}

\paragraph{Lean proof endpoint (built separately)}\begin{ReviewVerbatim}
GKGMM19IterativeLocalVoting.modelA_response_lp_normed_cost_minimizer_formula
\end{ReviewVerbatim}

\paragraph{Recorded source-to-Spec screening}
\textbf{Verdict:} \texttt{not recorded}
\\
\textbf{Reason:} not recorded
\reviewmatch{Matches source input}
\noindent\textbf{Reviewer annotation}\par
\reviewerbox
\clearpage
\hypertarget{source-claim-a9ebc621f1ff8c32}{}
\section*{Review row 31}
\textbf{Source locator:} {\footnotesize\raggedright cited publication, lines 118-\allowbreak{}145\par}
\paragraph{Verbatim source input}
\begin{ReviewVerbatim}
bounds on the sum of absolute values of the changes, the sum of the square of the changes,
and the maximum change, respectively.)
More formally, consider a M-dimensional societal decision problem in X ⊂RM and a
population of voters V, where each voter v ∈V has bounded utility fv(x) ∈R, ∀x ∈X.
Then we consider the class of algorithms described in Algorithm 1. We study the algorithm
class under two plausible models of how voters respond to query (1), which asks for the
voter's favorite point in a local region.

• Model A:
One possibility is that voters exactly perform the maximization asked of
them, responding with their favorite point in the given Lq norm constraint set. In other
words, they return a point arg maxx∈{s:||s−xt−1||q≤rt}fvt(x). Note that by de[U+001C control character]nition of this
movement, the algorithm is myopically incentive compatible: if a voter is the last voter
and no projections are used, then truthfully performing this movement is the dominant
strategy. In general, the mechanism is not globally incentive compatible, nor incentive
compatible with projections onto the feasible region. Simple examples of manipulations
in both instances exist.

• Model B: On the other hand, voters may not actually search within the constraint set
to [U+001C control character]nd their favorite point inside of it. Rather, a voter v may have an idea about how
to best improve the current point and then move in that direction to the boundary of
the given constraint set. This model leads to a voter moving the current solution in the
direction of the gradient of her utility function, returning a point xt−1 + rt
gt
||gt||q , for some
gt ∈∂fvt(xt−1). Note that ∂f(x) denotes the set of subgradients of a function f at x, i.e.
g ∈∂f(x) if ∀y, f(y) −f(x) ≥gT (y −x).
\end{ReviewVerbatim}


\paragraph{Expanded PaperInterface specification}
\begin{ReviewVerbatim}
∀ {Coord : Type u_1} [inst : Fintype Coord] [inst_1 : Nonempty Coord] (q : GKGMM19IterativeLocalVoting.SourceNorm)
  (center : Coord → ℝ) (r : ℝ) (gradient response : Coord → ℝ),
  GKGMM19IterativeLocalVoting.ModelBFiniteResponseAt q center r gradient response ↔
    response = fun i => center i + r * (gradient i / GKGMM19IterativeLocalVoting.finiteCoordinateNorm q gradient)
\end{ReviewVerbatim}

\paragraph{Lean proof endpoint (built separately)}\begin{ReviewVerbatim}
GKGMM19IterativeLocalVoting.modelB_finite_response_formula
\end{ReviewVerbatim}

\paragraph{Recorded source-to-Spec screening}
\textbf{Verdict:} \texttt{not recorded}
\\
\textbf{Reason:} not recorded
\reviewmatch{Matches source input}
\noindent\textbf{Reviewer annotation}\par
\reviewerbox
\clearpage
\hypertarget{source-claim-057affe0833911ec}{}
\section*{Review row 32}
\textbf{Source locator:} {\footnotesize\raggedright cited publication, lines 118-\allowbreak{}145\par}
\paragraph{Verbatim source input}
\begin{ReviewVerbatim}
bounds on the sum of absolute values of the changes, the sum of the square of the changes,
and the maximum change, respectively.)
More formally, consider a M-dimensional societal decision problem in X ⊂RM and a
population of voters V, where each voter v ∈V has bounded utility fv(x) ∈R, ∀x ∈X.
Then we consider the class of algorithms described in Algorithm 1. We study the algorithm
class under two plausible models of how voters respond to query (1), which asks for the
voter's favorite point in a local region.

• Model A:
One possibility is that voters exactly perform the maximization asked of
them, responding with their favorite point in the given Lq norm constraint set. In other
words, they return a point arg maxx∈{s:||s−xt−1||q≤rt}fvt(x). Note that by de[U+001C control character]nition of this
movement, the algorithm is myopically incentive compatible: if a voter is the last voter
and no projections are used, then truthfully performing this movement is the dominant
strategy. In general, the mechanism is not globally incentive compatible, nor incentive
compatible with projections onto the feasible region. Simple examples of manipulations
in both instances exist.

• Model B: On the other hand, voters may not actually search within the constraint set
to [U+001C control character]nd their favorite point inside of it. Rather, a voter v may have an idea about how
to best improve the current point and then move in that direction to the boundary of
the given constraint set. This model leads to a voter moving the current solution in the
direction of the gradient of her utility function, returning a point xt−1 + rt
gt
||gt||q , for some
gt ∈∂fvt(xt−1). Note that ∂f(x) denotes the set of subgradients of a function f at x, i.e.
g ∈∂f(x) if ∀y, f(y) −f(x) ≥gT (y −x).
\end{ReviewVerbatim}


\paragraph{Expanded PaperInterface specification}
\begin{ReviewVerbatim}
∀ {Coord : Type u_1} [inst : Fintype Coord] [inst_1 : Nonempty Coord] {p q : ℝ},
  GKGMM19IterativeLocalVoting.HolderDualFinite p q →
    ∀ {center ideal : Coord → ℝ},
      (∀ (i : Coord), center i ≠ ideal i) →
        ∀ (r : ℝ) (response : Coord → ℝ),
          GKGMM19IterativeLocalVoting.ModelBFiniteResponseAt (GKGMM19IterativeLocalVoting.SourceNorm.lp q) center r
              (fun i => -GKGMM19IterativeLocalVoting.lpCostGradientCandidate p (fun j => center j - ideal j) i)
              response ↔
            response = fun i =>
              center i - r * GKGMM19IterativeLocalVoting.lpCostGradientCandidate p (fun j => center j - ideal j) i
\end{ReviewVerbatim}

\paragraph{Lean proof endpoint (built separately)}\begin{ReviewVerbatim}
GKGMM19IterativeLocalVoting.modelB_finite_response_neg_lp_cost_gradient_formula
\end{ReviewVerbatim}

\paragraph{Recorded source-to-Spec screening}
\textbf{Verdict:} \texttt{not recorded}
\\
\textbf{Reason:} not recorded
\reviewmatch{Matches source input}
\noindent\textbf{Reviewer annotation}\par
\reviewerbox
\clearpage
\hypertarget{source-claim-0bfa667fe6808e39}{}
\section*{Review row 33}
\textbf{Source locator:} {\footnotesize\raggedright cited publication, lines 480-\allowbreak{}535\par}
\paragraph{Verbatim source input}
\begin{ReviewVerbatim}
sampling such a voter decreases fast enough).
Next, we introduce another general class of utility functions, which we call Weighted
Euclidean utilities, for which one can obtain convergence to a unique solution.

De[U+001C control character]nition 2. Weighted Euclidean utilities. Let the solution space X be decomposable
into K di[U+001B control character]erent sub-spaces, so that x = (x1, . . . , xK) for each x ∈X (where �K
k=1 dim(xk) =
M). Suppose that the utility function of the voter v is

fv(x) = −

K
�

k=1

wk
v
||wv||2
||xk −xk
v||2.

where wv is a voter-speci[U+001C control character]c weight vector, then the function is a Weighted Euclidean utility
function. We further assume that wv ∈W ⊂RK
+ and xv are independently drawn for each
voter v from a joint probability distribution with a bounded and measurable density function,
with W nonempty, bounded, closed, and convex.

This utility function can be interpreted as follows: the decision-making problem is de-
composable into K sub-problems, and each voter v has an ideal point xk
v and a weight wk
v
for each sub-problem k, so that the voter's disutility for a solution is the weighted sum of
the Euclidean distances to the ideal points in each sub-problems. Such utility functions
may emerge in facility location problems, for example, where voters have preferences on

323


[form-feed in source extraction]

Garg, Kamble, Goel, Marn, & Munagala

the locations of multiple facilities on a map. This utility form is also the one most closely
related to the existing literature on Directional Equilibria and Quadratic Voting, in which
preferences are linear. To recover the weighted linear preferences case, set K = M, with
each sub-space of dimension 1. In this case, the following holds:

Proposition 1. Suppose that conditions C1, C2, and C3 are satis[U+001C control character]ed, the voter utilities are
Weighted Euclidean, and voters correctly respond to query (1) according to either Model A
or Model B. Then, ILV with L2 neighborhoods converges with probability 1 to the societal
optimal point.

The intuition for the result is as follows: as long as the neighborhood does not contain
the ideal point of the sampled voter, the correct response to query (1) under weighted
Euclidean preferences is to move the solution in the direction of the ideal point to the
neighborhood boundary, which, as it turns out, is the same as the direction of the gradient.
Thus with radius rt, the e[U+001B control character]ective movement is
\end{ReviewVerbatim}


\paragraph{Expanded PaperInterface specification}
\begin{ReviewVerbatim}
∀ {Voter : Type u_1} {Coord : Type u_2} [inst : Fintype Coord] [inst_1 : Nonempty Coord]
  (E : GKGMM19IterativeLocalVoting.ILVEnvironment Voter (Coord → ℝ)),
  GKGMM19IterativeLocalVoting.proposition1_weighted_euclidean_l2 E = GKGMM19IterativeLocalVoting.proposition1Statement E
\end{ReviewVerbatim}

\paragraph{Lean proof endpoint (built separately)}\begin{ReviewVerbatim}
GKGMM19IterativeLocalVoting.proposition1_weighted_euclidean_l2
\end{ReviewVerbatim}

\paragraph{Recorded source-to-Spec screening}
\textbf{Verdict:} \texttt{not recorded}
\\
\textbf{Reason:} not recorded
\reviewmatch{Matches source input}
\noindent\textbf{Reviewer annotation}\par
\reviewerbox
\clearpage
\hypertarget{source-claim-732e45903404fd5e}{}
\section*{Review row 34}
\textbf{Source locator:} {\footnotesize\raggedright cited publication, lines 540-\allowbreak{}575\par}
\paragraph{Verbatim source input}
\begin{ReviewVerbatim}
normalized (i.e. not divided by ||w||2), the algorithm would converge to the same point, as
if utility functions were normalized.

3.2 Decomposable Utilities

Next consider the general class of decomposable utilities, motivated by the fact that the
algorithm with L∞neighborhoods is of special interest since they are easy for humans to
understand: one can change each dimension up to a certain amount, independent of the
others.

De[U+001C control character]nition 3. Decomposable utilities. A voter utility function is decomposable if there
exists concave functions fm
v
for m ∈{1 . . . M} such that fv(x) = �M
m=1 fm
v (xm).

If the utility functions for the voters are decomposable, then we can show that our
algorithm under L∞neighborhoods converges to the vector of medians of voters' ideal points
on each dimension. Suppose that hm
X is the marginal density function of the random variable
xm
v , and let ¯xm be the set of medians of xm
v . (By set of medians, we mean the set of points
such that, on each dimension, the mass of voters with ideal points above and below.)

Proposition 2. Suppose that conditions C1, C2, and C3 are satis[U+001C control character]ed, the voter utilities
are decomposable, and voters respond to query (1) according to either Model A or Model
B. Then, ILV with L∞neighborhoods converges with probability 1 to a point in the set of
medians ¯x.

Although simply eliciting each agent's optimal solution and computing the vector of
median allocations on each dimension is a viable approach in the case of decomposable
utilities, deciding an optimal allocation across multiple dimensions is a more challenging
cognitive task than deciding whether one wants to increase or decrease each dimension
relative to the current solution (see Section 5.2.3 for experimental evidence). In fact, in this
\end{ReviewVerbatim}


\paragraph{Expanded PaperInterface specification}
\begin{ReviewVerbatim}
∀ {Voter Coord : Type} [inst : Fintype Coord] [inst_1 : Nonempty Coord]
  (E : GKGMM19IterativeLocalVoting.ILVEnvironment Voter (Coord → ℝ)),
  GKGMM19IterativeLocalVoting.proposition2_decomposable_linf_medians E =
    GKGMM19IterativeLocalVoting.proposition2Statement E
\end{ReviewVerbatim}

\paragraph{Lean proof endpoint (built separately)}\begin{ReviewVerbatim}
GKGMM19IterativeLocalVoting.proposition2_decomposable_linf_medians
\end{ReviewVerbatim}

\paragraph{Recorded source-to-Spec screening}
\textbf{Verdict:} \texttt{not recorded}
\\
\textbf{Reason:} not recorded
\reviewmatch{Matches source input}
\noindent\textbf{Reviewer annotation}\par
\reviewerbox
\clearpage
\hypertarget{source-claim-4f41e995c4d500c9}{}
\section*{Review row 35}
\textbf{Source locator:} {\footnotesize\raggedright cited publication, lines 1944-\allowbreak{}1965\par}
\paragraph{Verbatim source input}
\begin{ReviewVerbatim}
Theorem 4. (Nemirovski et al., 2009; Strassen, 1965) Let θ ∈Θ be a random vector with
distribution P. Let ¯f(x) = E[f(x, θ)] =
�

Θ f(x, θ)dP(θ), for x ∈X, a non-empty bounded
closed convex set, and assume the expectation is well-de[U+001C control character]ned and [U+001C control character]nite valued. Suppose that
f(·, θ), θ ∈Θ is convex and ¯f(·) is continuous and [U+001C control character]nite valued in a neighborhood of point
x.
For each θ, choose any g(x, θ) ∈∂f(x, θ). Then, there exists ¯g(x) ∈∂¯f(x) s.t. ¯g(x) =
Eθ[g(x, θ)].

This theorem says that the expected value of the sub-gradient of the utility at any point
x across voters is a subgradient of the societal utility at x, irrespective of how the voters
choose the subgradient when there are multiple subgradients, i.e., when the utility function
is not di[U+001B control character]erentiable. This key result allows us to use the subgradient of utility function of a
sampled voter as an unbaised estimate of the societal subgradient.
Now, consider a convex function f on a non-empty bounded closed convex set X ⊂RM,
and use [·]X to designate the projection operator. Starting with some x0 ∈X, consider the
SSGM update rule xt = [xt−1 −rt(¯gt + zt + bt)]X , where zt is a zero-mean random variable
and bt is a constant, and ¯gt ∈∂f(xt). Let Et[·] be the conditional expectation given Ft, the
σ-[U+001C control character]eld generated by x0, x1, . . . , xt. Then we have the following convergence result.
\end{ReviewVerbatim}


\paragraph{Expanded PaperInterface specification}
\begin{ReviewVerbatim}
∀ {Theta : Type u_1} {Coord : Type u_2} [inst : MeasurableSpace Theta] [inst_1 : Fintype Coord]
  (mu : MeasureTheory.Measure Theta) [MeasureTheory.IsProbabilityMeasure mu] (solutionSpace : Set (Coord → ℝ))
  (sampleCost : Theta → (Coord → ℝ) → ℝ) (x : Coord → ℝ) (sampleGradient : Theta → Coord → ℝ),
  solutionSpace.Nonempty →
    Bornology.IsBounded solutionSpace →
      IsClosed solutionSpace →
        Convex ℝ solutionSpace →
          x ∈ solutionSpace →
            (∀ y ∈ solutionSpace, MeasureTheory.Integrable (fun theta => sampleCost theta y) mu) →
              (∀ (i : Coord), MeasureTheory.Integrable (fun theta => sampleGradient theta i) mu) →
                (∀ (theta : Theta), ConvexOn ℝ solutionSpace (sampleCost theta)) →
                  ContinuousAt (fun y => ∫ (theta : Theta), sampleCost theta y ∂mu) x →
                    (∀ (theta : Theta),
                        GKGMM19IterativeLocalVoting.FiniteSubgradientWithinAt (sampleCost theta) solutionSpace x
                          (sampleGradient theta)) →
                      GKGMM19IterativeLocalVoting.ExpectedSubgradientTheoremStatement mu solutionSpace sampleCost x
                        sampleGradient
\end{ReviewVerbatim}

\paragraph{Lean proof endpoint (built separately)}\begin{ReviewVerbatim}
GKGMM19IterativeLocalVoting.appendix_theorem4_expected_subgradient_boundary
\end{ReviewVerbatim}

\paragraph{Recorded source-to-Spec screening}
\textbf{Verdict:} \texttt{not recorded}
\\
\textbf{Reason:} not recorded
\reviewmatch{Matches source input}
\noindent\textbf{Reviewer annotation}\par
\reviewerbox
\clearpage
\hypertarget{source-claim-0f5dabcd51f758b4}{}
\section*{Review row 36}
\textbf{Source locator:} {\footnotesize\raggedright cited publication, lines 1971-\allowbreak{}2005\par}
\paragraph{Verbatim source input}
\begin{ReviewVerbatim}
Theorem 5. (Jiang & Walrand, 2010) Consider the above update rule. If

f(·) has a unique minimizer x∗∈X

rt > 0,
�

t
rt = ∞,
�

t
r2
t < ∞

∃C1 ∈R < ∞s.t. ||∂f(x)||2 ≤C1, ∀x ∈X

∃C2 ∈R < ∞s.t. Et[||zt||2] ≤C2, ∀t

∃C3 ∈R < ∞s.t. ||bt||2 ≤C3, ∀t
�

t
rt||bt||< ∞w.p. 1

Then xt →x∗w.p. 1 as t →∞.

Note: Jiang and Walrand (2010) prove the result for gradients, though the same proof
follows for subgradients. Only the inequality [x∗−xt]T gt ≤f(x∗) −f(xt) for gradient gt
at iteration t is used, which holds for subgradients. Boyd and Mutapcic (2006) provide
a general discussion of subgradient methods, along with similar results. Shor (1998), in
Theorem 46, provide a convergence proof for the stochastic subgradient method without
projections and the extra noise terms.

C.2 Mapping ILV to SSGM
\end{ReviewVerbatim}


\paragraph{Expanded PaperInterface specification}
\begin{ReviewVerbatim}
∀ {Omega : Type u_1} {Coord : Type u_2} {mOmega : MeasurableSpace Omega} [inst : Fintype Coord]
  [inst_1 : Nonempty Coord] (mu : MeasureTheory.Measure Omega) [inst_2 : MeasureTheory.IsProbabilityMeasure mu]
  (filtration : MeasureTheory.Filtration ℕ mOmega) (solutionSpace : Set (Coord → ℝ)) (objective : (Coord → ℝ) → ℝ)
  (project : (Coord → ℝ) → Coord → ℝ) (trajectory meanSubgradient noise : ℕ → Omega → Coord → ℝ) (bias : ℕ → Coord → ℝ)
  (radius : ℕ → ℝ) (xstar : Coord → ℝ),
  GKGMM19IterativeLocalVoting.AppendixTheorem5Statement mu filtration solutionSpace objective project trajectory
    meanSubgradient noise bias radius xstar
\end{ReviewVerbatim}

\paragraph{Lean proof endpoint (built separately)}\begin{ReviewVerbatim}
GKGMM19IterativeLocalVoting.appendix_theorem5_ssgm_convergence_boundary
\end{ReviewVerbatim}

\paragraph{Recorded source-to-Spec screening}
\textbf{Verdict:} \texttt{not recorded}
\\
\textbf{Reason:} not recorded
\reviewmatch{Matches source input}
\noindent\textbf{Reviewer annotation}\par
\reviewerbox
\clearpage
\hypertarget{source-claim-6ed3807a13e71f2b}{}
\section*{Review row 37}
\textbf{Source locator:} {\footnotesize\raggedright cited publication, lines 645-\allowbreak{}680\par}
\paragraph{Verbatim source input}
\begin{ReviewVerbatim}
325


[form-feed in source extraction]

Garg, Kamble, Goel, Marn, & Munagala

Rather, we consider an unconstrained budget allocation problem, one in which a voter's
utility includes a term for the budget de[U+001C control character]cit. Let E ⊆{1 . . . M}, I = {1 . . . M} \ E be
the expenditure and income items, respectively. Then the general budget utility function is
fv(x) = gv(x) −d(�
e∈E xe −�
i∈I xi), where d is an increasing function on the de[U+001C control character]cit.
For example, suppose a voter's disutility was proportional to the square of the budget
de[U+001C control character]cit (she especially dislikes large budget de[U+001C control character]cits); then, this term adds complex depen-
dencies between the budget items. In general, nothing is known about convergence of Algo-
rithm 1 with such utilities, as the de[U+001C control character]cit term may add complex dependencies between the
dimensions. However, if the voter utility functions are decomposable across the dimensions
and L∞neighborhoods used, then the results of Section 3.2 can be applied. We propose
the following class of decomposable utility functions for the budgeting problem, achieved
by assuming that the cost for the de[U+001C control character]cit is linear, and call the class [U+0010 control character]decomposable with a
linear cost for de[U+001C control character]cit," or DLCD.

De[U+001C control character]nition 4. Let fv(x) be DLCD if

fv(x) =

M
�

m=1
fm
v (xm) −wv

��

e∈E
\end{ReviewVerbatim}


\paragraph{Expanded PaperInterface specification}
\begin{ReviewVerbatim}
∀ {Dim : Type u_1} [inst : Fintype Dim] (utility : (Dim → ℝ) → ℝ) (isExpense : Dim → Bool)
  (componentUtility : Dim → ℝ → ℝ) (deficitWeight : ℝ),
  GKGMM19IterativeLocalVoting.paper_definition4_dlcd_formula utility isExpense componentUtility deficitWeight =
    (0 ≤ deficitWeight ∧
      (∀ (m : Dim), ConcaveOn ℝ Set.univ (componentUtility m)) ∧
        ∀ (x : Dim → ℝ),
          utility x = GKGMM19IterativeLocalVoting.dlcdBudgetUtility isExpense componentUtility deficitWeight x)
\end{ReviewVerbatim}

\paragraph{Lean proof endpoint (built separately)}\begin{ReviewVerbatim}
GKGMM19IterativeLocalVoting.paper_definition4_dlcd_formula
\end{ReviewVerbatim}

\paragraph{Recorded source-to-Spec screening}
\textbf{Verdict:} \texttt{not recorded}
\\
\textbf{Reason:} not recorded
\reviewmatch{Matches source input}
\noindent\textbf{Reviewer annotation}\par
\reviewerbox
\clearpage
\hypertarget{source-claim-9b6fb18f85b1c2c8}{}
\section*{Review row 38}
\textbf{Source locator:} {\footnotesize\raggedright cited publication, lines 645-\allowbreak{}680\par}
\paragraph{Verbatim source input}
\begin{ReviewVerbatim}
325


[form-feed in source extraction]

Garg, Kamble, Goel, Marn, & Munagala

Rather, we consider an unconstrained budget allocation problem, one in which a voter's
utility includes a term for the budget de[U+001C control character]cit. Let E ⊆{1 . . . M}, I = {1 . . . M} \ E be
the expenditure and income items, respectively. Then the general budget utility function is
fv(x) = gv(x) −d(�
e∈E xe −�
i∈I xi), where d is an increasing function on the de[U+001C control character]cit.
For example, suppose a voter's disutility was proportional to the square of the budget
de[U+001C control character]cit (she especially dislikes large budget de[U+001C control character]cits); then, this term adds complex depen-
dencies between the budget items. In general, nothing is known about convergence of Algo-
rithm 1 with such utilities, as the de[U+001C control character]cit term may add complex dependencies between the
dimensions. However, if the voter utility functions are decomposable across the dimensions
and L∞neighborhoods used, then the results of Section 3.2 can be applied. We propose
the following class of decomposable utility functions for the budgeting problem, achieved
by assuming that the cost for the de[U+001C control character]cit is linear, and call the class [U+0010 control character]decomposable with a
linear cost for de[U+001C control character]cit," or DLCD.

De[U+001C control character]nition 4. Let fv(x) be DLCD if

fv(x) =

M
�

m=1
fm
v (xm) −wv

��

e∈E
\end{ReviewVerbatim}


\paragraph{Expanded PaperInterface specification}
\begin{ReviewVerbatim}
∀ {Dim : Type u_1} [inst : Fintype Dim] (isExpense : Dim → Bool) (componentUtility : Dim → ℝ → ℝ) (deficitWeight : ℝ)
  (x : Dim → ℝ),
  GKGMM19IterativeLocalVoting.dlcdBudgetUtility isExpense componentUtility deficitWeight x =
    ∑ m, componentUtility m (x m) -
      deficitWeight * ((∑ m, if isExpense m = true then x m else 0) - ∑ m, if isExpense m = true then 0 else x m)
\end{ReviewVerbatim}

\paragraph{Lean proof endpoint (built separately)}\begin{ReviewVerbatim}
GKGMM19IterativeLocalVoting.dlcdBudgetUtility
\end{ReviewVerbatim}

\paragraph{Recorded source-to-Spec screening}
\textbf{Verdict:} \texttt{not recorded}
\\
\textbf{Reason:} not recorded
\reviewmatch{Matches source input}
\noindent\textbf{Reviewer annotation}\par
\reviewerbox
\clearpage
\hypertarget{source-claim-0c5fe9d4037e07b4}{}
\section*{Review row 39}
\textbf{Source locator:} {\footnotesize\raggedright cited publication, lines 430-\allowbreak{}450\par}
\paragraph{Verbatim source input}
\begin{ReviewVerbatim}
3.1 Spatial Utilities

Here we consider spatial utility functions, where the utilities of each voters can be expressed
in the form of some kind of spatial distance from their ideal solutions. First, we consider
the following kind of utilities.

De[U+001C control character]nition 1. Lp normed utilities. The voter utility function is Lp normed if fv(x) =
−||x −xv||p, ∀x ∈X.

Under such utilities, for p = 1, 2 and ∞, restricting voters to a ball in the dual norm
leads to convergence to the societal optimum.

Theorem 1. Suppose that conditions C1, C2, and C3 are satis[U+001C control character]ed, the voter utilities are Lp

normed, and voters respond to query (1) according to either Model A or Model B. Then,
ILV with Lq neighborhoods converges to the societal optimal point w.p. 1 when (p, q) = (2, 2),
(1, ∞), or (∞, 1).

322
\end{ReviewVerbatim}


\paragraph{Expanded PaperInterface specification}
\begin{ReviewVerbatim}
∀ {Voter : Type u_1} {Coord : Type u_2} [inst : Fintype Coord] [inst_1 : Nonempty Coord]
  (E : GKGMM19IterativeLocalVoting.ILVEnvironment Voter (Coord → ℝ)),
  GKGMM19IterativeLocalVoting.theorem1_lp_normed_dual_cases E = GKGMM19IterativeLocalVoting.theorem1Statement E
\end{ReviewVerbatim}

\paragraph{Lean proof endpoint (built separately)}\begin{ReviewVerbatim}
GKGMM19IterativeLocalVoting.theorem1_lp_normed_dual_cases
\end{ReviewVerbatim}

\paragraph{Recorded source-to-Spec screening}
\textbf{Verdict:} \texttt{not recorded}
\\
\textbf{Reason:} not recorded
\reviewmatch{Matches source input}
\noindent\textbf{Reviewer annotation}\par
\reviewerbox
\clearpage
\hypertarget{source-claim-ef1266fd9eeea090}{}
\section*{Review row 40}
\textbf{Source locator:} {\footnotesize\raggedright cited publication, lines 430-\allowbreak{}450\par}
\paragraph{Verbatim source input}
\begin{ReviewVerbatim}
3.1 Spatial Utilities

Here we consider spatial utility functions, where the utilities of each voters can be expressed
in the form of some kind of spatial distance from their ideal solutions. First, we consider
the following kind of utilities.

De[U+001C control character]nition 1. Lp normed utilities. The voter utility function is Lp normed if fv(x) =
−||x −xv||p, ∀x ∈X.

Under such utilities, for p = 1, 2 and ∞, restricting voters to a ball in the dual norm
leads to convergence to the societal optimum.

Theorem 1. Suppose that conditions C1, C2, and C3 are satis[U+001C control character]ed, the voter utilities are Lp

normed, and voters respond to query (1) according to either Model A or Model B. Then,
ILV with Lq neighborhoods converges to the societal optimal point w.p. 1 when (p, q) = (2, 2),
(1, ∞), or (∞, 1).

322
\end{ReviewVerbatim}


\paragraph{Expanded PaperInterface specification}
\begin{ReviewVerbatim}
GKGMM19IterativeLocalVoting.Theorem1NormPair GKGMM19IterativeLocalVoting.SourceNorm.l1
  GKGMM19IterativeLocalVoting.SourceNorm.linfty
\end{ReviewVerbatim}

\paragraph{Lean proof endpoint (built separately)}\begin{ReviewVerbatim}
GKGMM19IterativeLocalVoting.theorem1_norm_pair_l1_linf
\end{ReviewVerbatim}

\paragraph{Recorded source-to-Spec screening}
\textbf{Verdict:} \texttt{not recorded}
\\
\textbf{Reason:} not recorded
\reviewmatch{Matches source input}
\noindent\textbf{Reviewer annotation}\par
\reviewerbox
\clearpage
\hypertarget{source-claim-1ffdcf0343191a12}{}
\section*{Review row 41}
\textbf{Source locator:} {\footnotesize\raggedright cited publication, lines 430-\allowbreak{}450\par}
\paragraph{Verbatim source input}
\begin{ReviewVerbatim}
3.1 Spatial Utilities

Here we consider spatial utility functions, where the utilities of each voters can be expressed
in the form of some kind of spatial distance from their ideal solutions. First, we consider
the following kind of utilities.

De[U+001C control character]nition 1. Lp normed utilities. The voter utility function is Lp normed if fv(x) =
−||x −xv||p, ∀x ∈X.

Under such utilities, for p = 1, 2 and ∞, restricting voters to a ball in the dual norm
leads to convergence to the societal optimum.

Theorem 1. Suppose that conditions C1, C2, and C3 are satis[U+001C control character]ed, the voter utilities are Lp

normed, and voters respond to query (1) according to either Model A or Model B. Then,
ILV with Lq neighborhoods converges to the societal optimal point w.p. 1 when (p, q) = (2, 2),
(1, ∞), or (∞, 1).

322
\end{ReviewVerbatim}


\paragraph{Expanded PaperInterface specification}
\begin{ReviewVerbatim}
GKGMM19IterativeLocalVoting.Theorem1NormPair GKGMM19IterativeLocalVoting.SourceNorm.l2
  GKGMM19IterativeLocalVoting.SourceNorm.l2
\end{ReviewVerbatim}

\paragraph{Lean proof endpoint (built separately)}\begin{ReviewVerbatim}
GKGMM19IterativeLocalVoting.theorem1_norm_pair_l2_l2
\end{ReviewVerbatim}

\paragraph{Recorded source-to-Spec screening}
\textbf{Verdict:} \texttt{not recorded}
\\
\textbf{Reason:} not recorded
\reviewmatch{Matches source input}
\noindent\textbf{Reviewer annotation}\par
\reviewerbox
\clearpage
\hypertarget{source-claim-eab074ea3e4eb48e}{}
\section*{Review row 42}
\textbf{Source locator:} {\footnotesize\raggedright cited publication, lines 430-\allowbreak{}450\par}
\paragraph{Verbatim source input}
\begin{ReviewVerbatim}
3.1 Spatial Utilities

Here we consider spatial utility functions, where the utilities of each voters can be expressed
in the form of some kind of spatial distance from their ideal solutions. First, we consider
the following kind of utilities.

De[U+001C control character]nition 1. Lp normed utilities. The voter utility function is Lp normed if fv(x) =
−||x −xv||p, ∀x ∈X.

Under such utilities, for p = 1, 2 and ∞, restricting voters to a ball in the dual norm
leads to convergence to the societal optimum.

Theorem 1. Suppose that conditions C1, C2, and C3 are satis[U+001C control character]ed, the voter utilities are Lp

normed, and voters respond to query (1) according to either Model A or Model B. Then,
ILV with Lq neighborhoods converges to the societal optimal point w.p. 1 when (p, q) = (2, 2),
(1, ∞), or (∞, 1).

322
\end{ReviewVerbatim}


\paragraph{Expanded PaperInterface specification}
\begin{ReviewVerbatim}
GKGMM19IterativeLocalVoting.Theorem1NormPair GKGMM19IterativeLocalVoting.SourceNorm.linfty
  GKGMM19IterativeLocalVoting.SourceNorm.l1
\end{ReviewVerbatim}

\paragraph{Lean proof endpoint (built separately)}\begin{ReviewVerbatim}
GKGMM19IterativeLocalVoting.theorem1_norm_pair_linf_l1
\end{ReviewVerbatim}

\paragraph{Recorded source-to-Spec screening}
\textbf{Verdict:} \texttt{not recorded}
\\
\textbf{Reason:} not recorded
\reviewmatch{Matches source input}
\noindent\textbf{Reviewer annotation}\par
\reviewerbox
\clearpage
\hypertarget{source-claim-8fdb269066d0c098}{}
\section*{Review row 43}
\textbf{Source locator:} {\footnotesize\raggedright cited publication, lines 460-\allowbreak{}480\par}
\paragraph{Verbatim source input}
\begin{ReviewVerbatim}
rule is not equivalent to SSGM in general, and we leave the convergence to the societal
optimum for general (p, q) as an open question.
Further, note that even if each voter could scale their utility function arbitrarily, the
algorithm would converge to the same point.
However, the general result does hold for general dual norms (p, q) if one assumes the
alternative behavior model.

Theorem 2. Suppose that conditions C1, C2, and C3 are satis[U+001C control character]ed, the voter utilities are
Lp normed, and voters respond to query (1) according to Model B. Then, ILV with Lq

neighborhoods converges to the societal optimal point w.p. 1 for any p > 0 and q > 0 such
that 1/p + 1/q = 1.

The proof is contained in the appendix. It uses the following property of Lp normed
utilities: the Lq norm of the gradient of these utilities at any point other than the ideal
point is constant. This fact, along with the voter behavior model, allows the algorithm to
implicitly capture the magnitude of the gradient of the utilities, and thus a direct mapping
to SSGM is obtained.
Note that the above result holds even if we assume that a voter moves to her ideal point
xv in case it falls within the neighborhood (since, as explained earlier, the probability of
sampling such a voter decreases fast enough).
\end{ReviewVerbatim}


\paragraph{Expanded PaperInterface specification}
\begin{ReviewVerbatim}
∀ {Voter : Type u_1} {Coord : Type u_2} [inst : Fintype Coord] [inst_1 : Nonempty Coord]
  (E : GKGMM19IterativeLocalVoting.ILVEnvironment Voter (Coord → ℝ)),
  GKGMM19IterativeLocalVoting.theorem2_modelB_holder_dual_norms E = GKGMM19IterativeLocalVoting.theorem2Statement E
\end{ReviewVerbatim}

\paragraph{Lean proof endpoint (built separately)}\begin{ReviewVerbatim}
GKGMM19IterativeLocalVoting.theorem2_modelB_holder_dual_norms
\end{ReviewVerbatim}

\paragraph{Recorded source-to-Spec screening}
\textbf{Verdict:} \texttt{not recorded}
\\
\textbf{Reason:} not recorded
\reviewmatch{Matches source input}
\noindent\textbf{Reviewer annotation}\par
\reviewerbox
\clearpage
\hypertarget{source-claim-970bf26e54daa2aa}{}
\section*{Review row 44}
\textbf{Source locator:} {\footnotesize\raggedright cited publication, lines 590-\allowbreak{}620\par}
\paragraph{Verbatim source input}
\begin{ReviewVerbatim}
ited to simulations. We show that with the radius decreasing as O(1

t ), the algorithm indeed
[U+001C control character]nds directional equilibria in the following sense: if under a few conditions a trajectory of
the algorithm converges to a point, then that point is a directional equilibrium.

Theorem 3. Suppose that C1, C2, and C3 are satis[U+001C control character]ed, and let G(x) ≜Ev
�
∇fv(x)
||∇fv(x)||2

�
.

Suppose, G(x) is uniformly continuous, L2 movement norm constraints are used, and voters
move according to Model B. If a trajectory {x}∞
t=1 of the algorithm converges to x∗, i.e.
xt →x∗, then x∗is a directional equilibrium, i.e. G(x∗) = 0.

The proof is in the appendix. It relies heavily on the continuity assumption: if a point x
is not a directional equilibrium, then the algorithm with step sizes O(1

t ) will with probability
1 leave any small region surrounding x: the net drift of the voter movements is away from
the region. We note that the necessary assumptions hold for all utility functions for which
convergence holds, using the L2 norm algorithm (e.g. weighted Euclidean utilities). It is
further possible to characterize other utility functions for which the equivalence holds: with
appropriate conditions on the distribution of voters and how f di[U+001B control character]ers among voters, the
conditions on G can be met.
We further conjecture that even under voter Model A, if Algorithm 2 converges, the [U+001C control character]xed
point is a Directional Equilibrium. Note that as rt →0, fv(y) can be linearly approximated
\end{ReviewVerbatim}


\paragraph{Expanded PaperInterface specification}
\begin{ReviewVerbatim}
∀ {Voter : Type u_1} {Coord : Type u_2} [inst : Fintype Voter] [inst_1 : MeasurableSpace Voter]
  [inst_2 : MeasurableSingletonClass Voter] [inst_3 : Fintype Coord] [inst_4 : Nonempty Coord]
  (E : GKGMM19IterativeLocalVoting.ILVEnvironment Voter (Coord → ℝ))
  (M : GKGMM19IterativeLocalVoting.FiniteCoordinateTheorem3SampledProjectedSourceSemantics E),
  E.solutionSpace = Set.univ → GKGMM19IterativeLocalVoting.theorem3Statement E
\end{ReviewVerbatim}

\paragraph{Lean proof endpoint (built separately)}\begin{ReviewVerbatim}
GKGMM19IterativeLocalVoting.theorem3_statement_of_full_sampled_projected_source_semantics_univ
\end{ReviewVerbatim}

\paragraph{Recorded source-to-Spec screening}
\textbf{Verdict:} \texttt{not recorded}
\\
\textbf{Reason:} not recorded
\reviewmatch{Matches source input}
\noindent\textbf{Reviewer annotation}\par
\reviewerbox

\end{document}
